% =====================================================================
%  R. Nielsen, ON THE CONDITIONS FOR A STRONG WILL.  Six Lectures.  Copenhagen: Gyldendal (F. Hegel), 1874.  83 pp.
%  (Om Betingelserne for en kraftig Villie.  Sex Forelæsninger.)
%  ENGLISH TRANSLATION, from transcription.tex (Danish).
%  GENERATED by ebook-method/betingelserne-villie/tr_build.py --write from this template and the four translators' parts
%  in ebook-method/betingelserne-villie/tr/ (TRANSLATOR-BRIEF.md; terms: GLOSSARY.md, the title and drive-will /
%  reason-will / spirit-will approved by Hans 2026-10-12).  Conventions: TRANSLATION-PLAYBOOK.md.  Page numbers of the
%  1874 print in the margin.  Emphasis as in the Danish: \emph = letterspacing, \textbf = bold.
% =====================================================================
\documentclass[12pt,a4paper,openany]{book}
\usepackage[utf8]{inputenc}
\usepackage[T1]{fontenc}
\usepackage{libertinus}
\usepackage{libertinust1math}
\usepackage[english]{babel}
\usepackage[protrusion=true,expansion=true]{microtype}
\usepackage{setspace}
\usepackage[margin=1.3in]{geometry}
\usepackage{fancyhdr}
\usepackage{hyperref}
\hypersetup{pdftitle={On the Conditions for a Strong Will}, pdfauthor={Rasmus Nielsen}, hidelinks}
\setlength{\headheight}{15pt}
\pagestyle{fancy}
\fancyhf{}
\fancyhead[LE,RO]{\thepage}
\fancyhead[C]{\textit{R. Nielsen: On the Conditions for a Strong Will (1874)}}
\setstretch{1.3}
\usepackage{marginnote}
\newcommand{\opage}[1]{\marginnote{\footnotesize\textit{[#1]}}}
% a lecture opens a new page: heading, a short rule (the print has a short wavy rule)
\newcommand{\lecture}[1]{\clearpage\begin{center}{\large #1}\\[0.4em]\rule{3em}{0.4pt}\end{center}\medskip\noindent}

\begin{document}

% PDF 11, the title page
\begin{titlepage}
\begin{center}
\vspace*{3em}
{\large On}\\[0.6em]
{\Huge the Conditions}\\[0.6em]
{\large for}\\[0.6em]
{\Huge a Strong Will.}\\[2.5em]
Six Lectures\\[1em]
by\\[1em]
{\large R. Nielsen.}\\[5em]
Copenhagen.\\[0.5em]
Published by Gyldendal's Bookshop (F. Hegel).\\[0.5em]
{\small Printed by J. H. Schultz.}\\[0.5em]
{\small 1874.}\\[3em]
{\small Translated from the Danish}
\end{center}
\end{titlepage}

% ===== TEXT (pp. 1--83) =====
% --- p. 1 ---
\setcounter{footnote}{0}%
\lecture{First Lecture.}\opage{1}

There is in the mind of men an irrepressible unrest, a dim, ungovernable
drive, which hunts after freedom. What, then, is freedom? A personal
good, and hence more than an external condition. If freedom were merely
an external condition, a state that could be reached by withdrawing from
outward bonds and constraint, then the Gypsy would be freer than the citizen,
and the wild beasts of the field and the birds under heaven
freer than man. But freedom is essentially an inward
good; freedom is the independence of our individuality, the fundamental
stamp of our personality; freedom is our patent of nobility, our inner,
true, original self.

Independence rests on self-subsistence, self-subsistence
on self-determining, and self-determining is freedom: the more
self there is in a human being, the greater the freedom; the
smaller and weaker the self, the smaller the freedom.
The ground of freedom is will, for the will is the core of the self.
The self is what determines itself, but the capacity for self-determination
is precisely the will; the will is free: the stronger the will, the
greater the freedom, the weaker the will, the smaller the freedom.

Now if all men are originally predisposed to freedom,
then it is given herewith that they all without exception must
have a free will. But while we assume that all men
have a free will, we must nevertheless admit that we are far from all
possessing the freedom of the will in equal measure; some have a strong,
others a weak will: whence can that possibly come? Is
the will something innate or something self-acquired or perhaps
% --- p. 2 ---
\setcounter{footnote}{0}%
\opage{2}a gift of heaven, a gift of the Spirit, the eternally free? We
may well say that it is in part something innate, in part
something self-acquired, and yet in each and every respect a gift of heaven;
but then it depends on how this is more closely
to be understood.

Our first presupposition is that the will rests on
something innate. To will belongs natural predisposition. It
is a fact that individualities are born with so
conspicuous an independence that they proceed from the hand of Nature
as pure will. % print: Villie (misprint)
They are the born rulers; the stamp of the ruler
with its superiority is upon them, whether they
otherwise find themselves in a higher or a lower station. Those
who saw a Napoleon in his first youth met the
same glance, the fine, legible traits of destiny, the same closedness
and iron-firm will in the young pupil from Brienne as
later in the great Emperor. Others, by contrast, come into the world
with so weak a predisposition to independence and betray so
striking a lack of will that they seem destined to walk
in leading-strings throughout their whole life. Even if they are born in
purple and, so to speak, with a scepter in their hand, the
weakness of will is nevertheless so great that they are never themselves able
either to grasp or to hold fast any decisive resolution.
They live under a perpetual guardianship: favorites,
servants, counselors must take it upon themselves to will in their stead,
and they take it upon themselves gladly.

Instead of proceeding from the principle of all
men's equal title to the goods of freedom, one would be tempted,
from this point of view, to believe that Nature had
set up an insurmountable boundary between the free-born and the
unfree-born, had predestined some to independence,
others to dependence, some to be masters, others slaves.
As far as the innate is concerned, however, we dare not overlook
that the faculty of will by no means stands in the same relation as
other faculties of the soul, for example memory. When someone says: I
once had a very faithful memory, but have unfortunately
almost lost it in a severe illness, we pity him,
and our pity is without reservation. If someone, on the other hand,
% --- p. 3 ---
\setcounter{footnote}{0}%
\opage{3}were to step forward and say: I once had a very strong
will, but unhappily lost it in a fever, we pity
him too, no doubt, but yet with a peculiar reservation. --- One
can lose one's sight without losing one's hearing, one can to a
certain degree lose one's memory without losing one's understanding;
but we know of no illness in which a human being can lose
the will and for the rest recover his health.

Whoever assures us that he has all the faculties of his soul unimpaired
except the will, which unfortunately was lost in an ague,
cannot be right in his head. A strange
human being! He has knowledge, education, a brilliant memory
and, as it seems, good understanding too, can join
in the talk on many things, science, politics, art, business affairs,
without one's noticing anything; only when he comes
to the one point, the point about the will and the ague,
does he become odd. Such a thing can happen when we think
of mental illnesses. An elderly woman, who in her family
had always been regarded as sound and hale, suddenly came
to the conviction that she could not eat. Had
she then no appetite? Yes. Was she afraid of being
poisoned? No. What then was the matter? Alas, it was
a hard fate. She had long suspected it; but now
it had suddenly dawned on her; she was, namely,
not, like other human beings, provided with inward organs,
and what use was it, then, that she ate? She was,
she assured us, hollow inside, quite hollow, hollow like a
drum, and then she wept. In a similar way a
human being can fall into the fixed idea that he has lost his will,
or has not been endowed with will at all, that in his
inner being there is something hollow, deep inside an empty space, a little
dark hole, just where the will should sit.

It is a fixed idea; but in feeling convinced
that it must be a fixed idea, we make an advance in
knowledge. For we see that the will, despite the innate
element in the faculty, is differently dependent on ourselves than any other
faculty whatsoever. I do not get a strong memory
because I will it, or a fruitful understanding, because I
% --- p. 4 ---
\setcounter{footnote}{0}%
\opage{4}will it, or a deep feeling or a rich imagination,
because I will it; such advantages rest on natural endowment,
and of these it holds that no one can add a cubit to his
stature. With the will it is another matter; will I get
only in that I will, solely and alone because I will. A human being
may have gifts and faculties ever so great: if
he does not will, he gets no will. He who wills
by halves gets a half will; but he who wills strongly gets
a strong will. From this point of view the will is its own
cause: by willing we produce will. The will is thus
our own production, a production in which the self,
the willing element in us, produces itself. It may be given to us to
will, but the will cannot be given to us, we must acquire it
ourselves: the will is something \emph{self-acquired}.

Why, then, do not all men have an equally strong
and vigorous will? According to the view last set forth, it must
be their own fault. And yet this view too has its deviation.
Here someone steps forward with
an objection; let us hear him. He says: If will
rested on willing alone, then my will, be it said
without any modesty, would be one of the greatest and
strongest of the age. For the self-acquisition of the will, will
produced by will, an ever increasing and strengthening
self-determining, a mighty, great and vigorous will is
precisely the ideal that I have been enthusiastic about from my youth.
But I have, just on account of my zeal, fallen short
of the ideal. I have not given myself over to the visionary
fancy that a human being can in the outward sphere do all that
he wills, that the will can move mountains. I have kept to
the faculties granted me; but precisely in the inward sphere I
have seen the yawning chasm between \emph{willing} and \emph{being able}.
In order to determine the will as my own, I had to
begin with a careful and exact separating of what
was mine from what was not mine. But when I
had separated off everything that was not mine, how much did I
have left? Only the willing, the inner determining,
the thought of will, pure willing, a shadow-being without power.
% --- p. 5 ---
\setcounter{footnote}{0}%
\opage{5}Only in so far as I am distinctly conscious of my will to myself, dare
I call it mine, and only in so far as I dare call it
mine, can I myself determine it. What is unconscious to me
is not mine; no one can determine what he does not
know. The only thing in the will of which I am conscious is this,
that I will. With this consciousness I can sit whole days
and will and will without accomplishing the least thing. The
moment I try to undertake anything, there at once
step in a multitude of heterogeneous attendant circumstances, which
I do not see through, and over which I am not master. I
may sit at a richly laid table, hungry and thirsty:
if I did nothing but will, I should starve
to death, and that was precisely % print: netcp (misprint)
what I did not will. Now
men will say, to be sure, that this is hair-splitting. ``Reach out your
hand and take for yourself; in other words: use your
powers, and then you will not starve to death.'' But I have a suspicion
that men in such matters do not always know
what they are themselves saying. ``Use your powers!'' that is precisely
my riddle. Is it the use of the powers that rests on the will,
or is it not rather the use of the will that rests on
the powers? When we voluntarily move arm and hand,
it seems to us, no doubt, that it is our conscious will that
does it; but that it is not so, the slightest
reflection can teach us. The will moves the nerves, the nerves the
muscles, and the muscles the arm, without our being in any
way able to become conscious how
it actually comes about. This much only we know, that it is not
that in our will which we dare call our own that is able
to accomplish such a thing. If the arm becomes lame, if the motor nerves are cut
across, then the person concerned may will as vigorously
as he wills: the arm does not move.

If the use of the powers depended unconditionally on our
will, then our will would have to have the powers in its
own power. To acquire will by strongly willing would then
be literally the same as to produce the bodily
powers which the will requires. One would then have to be able, by
self-acquired will, to overcome illness and produce
% --- p. 6 ---
\setcounter{footnote}{0}%
\opage{6}health; but that this is impossible everyone admits. It
is a fact that certain muscles and muscular movements
are subject to the will; but that it is by no means the will,
our conscious, self-acquired will, that from the beginning
has made itself master of the muscles, is obvious.
To want to explain the use of the organs by habit is thoughtless.
Habit comes by repetition; the will, it is said, repeats
the use of the organs until at last it
becomes master of the movements. So one explains the use
of the powers by saying: the use repeats itself; a poor
explanation when the very question is of the possibility of the use.
It is as if one wished to explain a miracle
by saying: the miracle repeats itself.

The illness that is called St. Vitus' dance gives
conspicuous examples of a disproportion between the will and
the movements. The limbs move without consulting
the will. --- In the stream of speech it seems as though the words in
their connection were directly governed by our conscious
will, and yet it is not always easy to find just the word
one seeks. In cases of weakening of the brain the divorce
between what the sick man actually wished to say and what
he really says can even be so great that he, for example,
says kitchen when he means bread, and drawer when he
means money. The will aims at the right thing, but the moving
powers elude the right use.

It is with the will as with a master who accomplishes
everything by servants, and for that reason must bear the responsibility
for all that the servants do. If it is something good, then
one praises, not the servants, but the master, for it was he
who willed it; if it is something evil, then one blames, not
the servants, but the master, for it was he who willed it.
This is all in order, so long as the servants are,
like will-less slaves, punctually obedient. But if they begin
to go their own way and become disobedient or even
only half obedient, then the master can no longer take the blame
upon himself. And yet a master has as a rule greater authority
over his servants than the will over its powers.
It
% --- p. 7 ---
\setcounter{footnote}{0}%
\opage{7}is within the power of a master to dismiss a disobedient servant
and take on a new one; but it is not within the power of the will to dismiss the old disobedient nerves and take on
new ones in their place.

Use your powers! it is as if an Aladdin were to
say: rub the lamp! The will's use of the powers is
not, to be sure, fabulous, for it is bound by laws; but in
its ultimate ground it is nevertheless just as inexplicable as
the coming of the spirit when the lamp is rubbed. The will wills, and
the powers stir; the lamp is rubbed, and the spirit
comes: the first phenomenon is as enigmatic as the
last.

We have here considered two mutually opposed ways of representing the matter,
of which the first makes the will exclusively
something \emph{innate}, while the last makes it just as
exclusively something \emph{self-acquired}. That they
both end in the absurd comes from this, that they both, each
in its own way, confuse the will itself with its conditions.
The innate is not yet will; it is only predisposition of will,
and so far a \emph{condition}. What is self-acquired is not the
ground of the will; self-acquisition presupposes precisely that which
is to be acquired. To produce will, one must will, but
in order to be \emph{able} to will, one must already have will; in order
to will strongly, one must have a strong will. He
who has no will at all to begin with cannot will,
and he who does not beforehand have a strong will cannot
will strongly. Thus it is seen from this that the self-acquisition
in which the will comes to be and grows cannot be the will's
innermost ground and first cause, but only a condition.


Now one might be tempted to throw the two conditions
together and say: the will is in part something innate, in part
something acquired; by taking the innate and the acquired
together, we get both the will's self-determining
and its ground, hence: the whole will. This is an impossibility;
the two conditions are far too heterogeneous for them
to be fused together in that way without more ado. There
% --- p. 8 ---
\setcounter{footnote}{0}%
\opage{8}is in the essence of the will something that rests on Nature,  and
something that rests on us. But what rests on Nature
cannot possibly rest on us; for if it did, we should have
to say that the whole rests on us. Conversely, what
rests on us cannot be said to rest on Nature either;
for if it did, the whole would have to
rest on Nature. No, what is to be able to unite
the two heterogeneous conditions must necessarily be a
third, a Higher, a power of will which is at once in Nature
and yet above Nature, in us and yet above us: this
power of will is \emph{spirit}.

When our free will is conceived as a gift of the Spirit, it becomes
in the highest sense a task, and the solution of the task
religious. But precisely with the religious the difficulties begin
anew, for here the misinterpretations begin in earnest.

The religious has many enemies. Pious men
complain bitterly of the ungodliness of the present and its increasing
unbelief; they may well have cause for this when they turn their gaze outward
and see the steadily growing number of external enemies. But
the enemy is within! ``Yes,'' they answer, ``our hereditary enemy
is within, it is the Devil, the Antichrist, the beast
with many horns, whose sign we shall soon see --- in
the external.'' But the enemy is within, in the religious
itself, an army of stealthy misunderstandings. The first
dangerous enemies that a believing mind has to defend itself against
are abortions of its own swelling fancy. Such
abortions show themselves under many disguises; they do not
always announce themselves with anguish of soul, shy timidity and doubt, they
announce themselves just as often with swollen faith, false security
and a certain thick-blooded self-confidence, which hinders the
development of the spirit. Those who would build upon religion's holy
ground must conduct themselves like the Israelites when they built
upon the second Temple, the sanctuary which in the fullness of time
was to receive the Savior of the world; they built with the weapon
in hand. Only the difference is this, that the enemies of that
outward Temple were outward; but the enemies of the sanctuary
of the spirit are essentially inward.

% --- p. 9 ---
\setcounter{footnote}{0}%
\opage{9}Reflections, too, are dangerous when the understanding, while
it is still uncertain about the presuppositions of the life of will,
begins to reason. Here is one who complains that in spite
of all his reflection he is not able to get clear about how things
actually stand in faith with regard to
spirit and will. I hear, to be sure, he says, a
good many unctuous assurances that it is God
the Holy Ghost who regenerates man, creates in us a
new man, a new mind, a new will, etc. When
this is spoken of thus in general, especially when it is
spoken of with pathos, I get a feeling as though I too
understood it, that is to say: in general. But when I
then go into myself and really try to give myself an account of
how things actually stand with spirit and will,
my head spins. The Awakened call out to me:
``You are not to ponder, but to believe; if you will only believe, it
goes of itself!'' I should like to believe, but with what
am I to begin practically, from whom is the beginning to proceed,
from the Spirit or from myself? If it is I who am to
take the first step, then I am compelled to begin
with the will I have; but the will I have is
precisely the old will, a will that is no good, a will
outside the Spirit, hence a spiritless will. How can
a human being come to the Spirit by means of a spiritless
will? Then it is said: ``You are to pray for the Spirit!'' But
as long as I have not yet received the Spirit, my prayer is
spiritless! So it cannot be I who am to begin;
it must be the Spirit. If it now stands firm that the beginning
proceeds, not from me, but from the Spirit: then I have nothing
else to do than to sit still and wait for the Spirit.
Am I then to become a Quaker? I am answered: ``You
have the Spirit, as surely as you are baptized. It depends solely
on yourself whether you will remember your baptismal covenant and
let the Spirit work.'' --- Let the Spirit work! But then
here is, after all, something that rests exclusively on me,
not on the Spirit. What does it mean to remember one's
baptismal covenant? It cannot possibly be a mere and bare
% --- p. 10 ---
\setcounter{footnote}{0}%
\opage{10}matter of memory; here there must be something that expressly
rests on mind and will, my own mind, my own
will. In order that the Spirit may be able to work in me, I must
therefore, as said above, by my still spiritless
will help it to work. So one
resorts to the last explanation and says: ``You are not to
spy on the Spirit, not to presume to will to understand how
it works. It works in secret as a new nature;
you are to experience its activity as a fact.
You can experience that the plant grows, but not see through
the secret of Nature and grasp how. The new life
is a new plant, it grows a quiet growth: we grow, we
sprout, that is a fact.''

So much edifying talk is then made about the spiritual
growth of life, how it sprouts even when we sleep. I cannot,
unfortunately, get sight of the fact. Barley sprouts
and oats sprout; it sprouts in the woman with child, and
the child is born and grows to a man, and there awakens in the man
a drive which the child did not know; those are facts.
But with the growth of life in believers it is a peculiar matter.
I have for many years heard them say that it sprouts in them,
that they are constantly growing; I have taken them at their word and
as exactly as possible taken their measure, so as at last
to be edified by the fact; but what have I
experienced? In the outward sphere I have indeed been able to see that the
believers too change with the years, that some of them
become, if not always taller, then at any rate not seldom a good deal
stouter; but as to the development of character I have not
been able to discover any progress. It seems
to me that the average stands exactly at the
same mark this year as ten, yes twenty years ago: decent
people who --- save for the turns of phrase --- have a
striking likeness to other decent people. This business of
growth must then be a great secret. My secret
is that it will not sprout in me. I too received an awakening
in my time and hoped that it would sprout;
% --- p. 11 ---
\setcounter{footnote}{0}%
\opage{11}but alas, it does not sprout: I am like the fir tree that
once had been a Christmas tree, --- set in growth.

This sample of reasoned-out misunderstandings contains
a forceful reminder that the enemy is in the
inward.

What, then, does all this aim at? An understanding % print: Fortsaaelse (misprint)
of what will is. We began by assuming the will to be
something innate; but the assumption led to the absurd.
We went on to conceive the will as something self-acquired,
but the conception led to the absurd. We ended by
referring the will to the Spirit; but as soon as the task became religious,
the misunderstandings came. This is therefore to introduce
an attempt to trace the will's origin and to \emph{clarify the conditions
for a strong will}.

% --- p. 12 ---
\setcounter{footnote}{0}%
\lecture{Second Lecture.}\opage{12}

According to an old division, man is said to consist
of three principal parts: \emph{body}, soul and \emph{spirit}. This tripartition
must not, however, be understood as if man % print: Mennnesket (misprint)
were composed of three pieces, and the whole had come about by
the soul's being put into the body, and the spirit into the soul,
so that the outer was to surround the inner as one box
can surround another. If the inner were merely boxed
within the outer, then unity would be lacking, the organic
unity that makes us living beings.

One of the strongest testimonies that the human
body, despite the multiplicity and diversity of its
parts, is not what is called a machine, a patched-together
whole, but is animated by life and unity of self, is the feeling
of pain. There can be pain in the most diverse parts
of the body, but pain as pain does not properly concern
the parts for themselves, not even the nerves; it is the soul
that feels. He who, to keep to everyday experiences,
has, for example, toothache does not say: that is the teeth's business, it
is a disturbance in the nerves which does not concern me; no, he
must confess that it concerns him, that it is he himself
who must feel the pain, the pain that is in the teeth.

Were man such a patched-together whole that
each part had its own life and its own feeling for itself,
then it would have to be easy for the fully conscious self to subject
its body to operations; but in order to subject itself
% --- p. 13 ---
\setcounter{footnote}{0}%
\opage{13}painful operations and yet, so far as possible, escape
the pain, one is nowadays chloroformed. By
this self-anesthetization it comes about that the separate parts of the body,
the separate organs, become like pieces that suffer each
for itself.

Feeling, the feeling of self especially, grows, is developed, is refined
with civilization; but it almost looks as though
the force of will were on the decline. When Mucius Scaevola thrust his hand
into the fire and let it burn, while he spoke to Porsenna,
he was certainly not chloroformed. It does indeed look
as though his hand were only a piece of body by itself,
as though it did not concern him, but were solely the hand's affair
how far it could endure the fire or not; and yet there was
no parcelling out. He felt what was going on, he felt
the pain; but there was in his self something that was
stronger than feeling, raised above the pain, and that was his
resolved will. Mucius Scaevola was a man of strong
will. Is it then perhaps only the man who, in relation to
pain, can have such a strong, heroic will? No!
the woman too. History knows how to speak of this in the life of
the Norsemen and of the Romans. When Marcus Brutus, Cassius and Casca
had conspired to murder Caesar and save the Republic,
Brutus bore the secret heavily and dared not
confide it to his wife: she was, after all, only a woman.
Portia suspected what was going on, but in order to make herself worthy
of an initiation and to show that she, a woman, could stand at
her husband's side even in a moment of horrors, she gave herself
a deep and painful stab in the foot. Then she stepped forward,
lightly and with a sure step, though the pain burned, and now
Brutus dared confide everything to her. Portia was a woman of
will, a great and firm and strong will.

It is thus clear from the will as well as from feeling that
our human being is not a patched-together whole,
but a unity of life, a self-determining unity of self, a
unity of feeling and will, an I.

Healthy common sense has in its own way raised itself above
the patching-together by lifting man up above the three parts of the
% --- p. 14 ---
\setcounter{footnote}{0}%
\opage{14}tripartition; every human being calls himself an I and says:
I have a body, I have a soul, I have spirit.

Here there enters a peculiar relation between
having and being. As regards our possession of external
things, it seems unnatural to make what we have
one with what we are. Because someone has an estate,
he does not say: I am what I have, an estate,
any more than he who has a walking-stick will say: I am a stick.
True, it is said that in America, in certain circles, it is the custom that
the host introduces two gentlemen in the company to each other briefly
and categorically as: two million dollars! --- four million! but,
although such a numerical designation has the advantage that it
exactly states the angle of inclination of the different nod with which
the two ought to meet each other: yet it cannot
be denied that so direct a fusion of having
with being must strike us as somewhat too materialistic.

It is otherwise when the talk is of the constituents of our being
and our essential attributes. In this connection
having can coincide in different ways with
being. We have a body, we are bodies, our being
is bodily; we have a soul, we are souls, our being is
soulish; we have spirit, we are spirits, our being is
spiritual. In this fusion of what we have and
what we are, we hold fast the unity of our being. % print: Eenhed (misprint)
Our self,
our I, the I we have and the I we are, is from the first
germ bodily-soulish-spiritual. There is in man
only one nature; but this nature has three principal sides, and
the whole nature is again in each of its principal sides. There
are in our fundamental nature three natures: a \emph{sensuous} one, corresponding
to bodiliness and that side of the soul which comes
under sensation and sensuous impulses; a \emph{rational} one,
corresponding to the incorporeal, hence to that side of the soul
which goes up into intuitions of ideas and the thinking of the universal;
a \emph{spiritual} one, corresponding to the original self in eternal freedom,
to the Spirit.

This tripartition is decisive for the will and its
conditions, for from it follows that the will too is threefold.
% --- p. 15 ---
\setcounter{footnote}{0}%
\opage{15}Our many troubles with freedom, all our exertions
to get ourselves in our own power and to work ourselves up to independence
stem from this, that we originally have three wills:
a \emph{drive-will}, corresponding to our sensuous nature and to what
is innate in the will, a \emph{reason-will}, corresponding to our
knowledge of the universally valid laws of the world and to what is
self-acquired in the will, and finally a \emph{spirit-will}, a will which,
when it understands itself and recognizes its origin,
has its independence in unconditional dependence on the Spirit
itself, and knows itself moved by God's will.

There are thus three stages on which man can stand
with his threefold will; but it is clear that we must stand
either on the first or on the second or on the
third of these stages, if we are to have an undivided, firm and
strong will.

\emph{The natural man}, by which we understand not merely
man as he is in a raw, natural condition, but man
immediately bound to his innate natural ground,
stands on the first stage. He is, with his individual
endowment, his reason and genius, as if cast in one piece.
The natural man has, regardless of how far his
civilization has reached, only one ruling will, one fundamental will,
and that is the \emph{drive-will}. The heroic life of the ancient North
and the flourishing cultural life among the Greeks teach us what
this means.

\emph{The man of reason}, by which we think exclusively
of such as have a will moved by what is necessary for thought
and universally valid in its essential necessity,
stands on the second stage. The thinker, the wise man, who only
wills what he must will, has in his agreement with
himself one will, one fundamental will, and that is: the \emph{reason-will}.
The great characters of antiquity, a Socrates, the Stoics
and others, furnish examples.

\emph{The man of spirit}, in the Archetype himself the only
true man, the Revealer of freedom, who wills to set all free,
stands exclusively on the third stage. Here there are in unity of will
two wills, a divine and a fully
% --- p. 16 ---
\setcounter{footnote}{0}%
\opage{16}human one; but the two wills are not parcelled out, not
divided; the divine is human and the
human divine; the unity of both is an absolutely
personal unity of will, and that is the \emph{spirit-will}.

But we, children of the nineteenth century, men of nature, of culture and
of spirit, of many broods, where do we now
actually stand? Not on any single one of the three stages; we stand
on all stages, all possible intermediate stages. That comes from this, that
we are so many-sided. To stand on one stage, to remain standing
and stand firm, is one-sided, and one-sidedness is the vice
that of all most must be feared. So we are many-sided, that
is to say: we have something of a drive-will, something of a reason-will
and a little of a spirit-will in various mixture.
This many-sidedness of ours means that the life of will, despite all its
bustle, goes somewhat cramped. We must serve not merely two,
but three masters; for each of the three wills in us is a
great power, and great powers are always domineering. If none
of them can become overlord, they stir up strife; they would weaken
us, parcel out our being and divide us, as the great powers in
their time divided Poland.

We shall in what follows consider each of the three
stages of will for itself, and begin with the \emph{drive-will}.

Man has drives like the animal, but for that reason
man is not an animal, and the human drives not
animal. One ought not to despise the drives, they belong necessarily
to our human nature; only when a human
drive begins to become animal has man reason
to become afraid of himself. The distinguishing mark of the human
drive is that the sensuous can take up the ideal,
come under the rational and suck strength from the spiritual.
The drive can embrace the lowest and the highest; there
is a drive to sensuous enjoyment, but there is also
a drive to the Kingdom of God.

The drive is the fruitful womb of the will, the sheath and
swaddling from which it unwinds itself. Therefore the drive-will works
with so wondrously primitive a natural force. The will cannot
grasp without being gripped by something, cannot move
% --- p. 17 ---
\setcounter{footnote}{0}%
\opage{17}with force without being inwardly and strongly moved by
something. That which from the beginning drives and moves,
because its object fills the will, the inner driving element of the natural willing,
is the drive.

The drive-will can appropriate the rational and be filled
with ideas, but it does not for that reason become reason-will; it
becomes only a \emph{rational} drive-will. It can have
the longings of blessedness, appropriate the promises of freedom, yes
even the spiritual essence of deep love, and yet it
never becomes spirit-will; it is and remains a spiritual drive-will.
Think of the relation between man and woman,
between mother and child, and one will recognize that the sensuous
can be transfigured into the ideal, the spiritual, on the basis
of the life of the drives. The passions can be at once
both sensuous and ideal. The poets praise love,
and we are grateful to them. But love is a
will, and yet one also calls it a passion. Right!
In Eros the will is passion, and the passion will;
so it must be: it is drive-will. Those who maintain
that this drive-will in its veiled ground is animal, slander
Eros and mock what is human. Juliet has a will
to Romeo, her will is passion; but this passion
is, as a spiritual atmosphere of life, the higher element in which she
grows; its glow and flame is the light in which she stands
transfigured.

How great the scope of the drive-will is can be seen when it is
applied to the determination of \emph{people} and \emph{multitude}. Where
is the people? In the multitude. Where, then, is the multitude? In % print: 1 (misprint)
the people.
And yet the multitude is not yet the people, and the people more
than the multitude, although both determinations fall within
the drive-will. We have a grand-scale example in the Old
Testament. When Moses was with Jehovah on the mountain, the
people broke loose, became a mob, became a multitude. This multitude, moved
by passion, demanded of Aaron that he should
make them gods who could lead them and ``go before
their face.'' That was drive-will. Gods of gold, golden gods,
that was something more real and present
% --- p. 18 ---
\setcounter{footnote}{0}%
\opage{18}than that Invisible One up there in the clouds. The multitude demanded
it; it was loud-voiced, for it had a will. Aaron
found that it was a strong will; so he confounded the
people with the multitude and cast the golden calf. When Moses
came, and Aaron saw his face, he grew afraid and prayed:
``let not my lord's wrath blaze up; thou knowest the people, that
it inclines to evil.'' Then it was shown what spirit-will
is. Moses placed himself at the gate of the camp and commanded: ``whoever
belongs to Jehovah, let him come to me!'' Moses staked
his life, for the mob was raging. But when the voice of command
sounded, the authoritatively commanding, the mob got other thoughts,
another will: the multitude became a people again. Did the people
then suddenly get a spirit-will? No! only Moses had spirit-will;
but the people got through him a will which from involuntary
drive had to bow before spirit; it was a spiritual drive-will,
and a higher will the people never gets as a people.

If we go into the particular forms, then every
intellectual striving that originally rests on special endowment,
talent, genius, is to be referred to the drive-will; will of genius is
drive-will.

It is a deep and true and touching trait in Oehlenschläger's
Correggio, that poor Antonio, whose courage
Michelangelo's harsh, unjust verdict has crushed, calls
himself to account and consoles himself with his artist's conscience,
saying: ``God knows, I did not do it
to shine. I painted as the bee builds cells, as
the bird builds its nest in a tree,'' that is, from inner urge,
from artist's drive; it is drive-will.

When it is a matter of practical outward-going activity, of
historical deeds, what significance do not drive-impulses
and passions have! We see it in the heroes, a
circle of great figures, all of whom stand under the drive-will.
How much that is rational, how many ideas and thoughts has
the will not here been able to appropriate; and yet it is,
when we look more closely, not reason-will, still less
spirit-will, but a drive-will inspirited by ideas, glowing through with passions,
that we have before us. It is a spiritual
% --- p. 19 ---
\setcounter{footnote}{0}%
\opage{19}natural force, an irresistible longing that drives the hero forward;
he wills because he must will. Earl Skule in ``The Pretenders''
cannot give up his struggle against Hakon, cannot renounce
the throne and make a settlement; and why not? His being's
urge is against it; the great royal thoughts say no.

It is a peculiar, uplifting sight to behold the many
highly different phenomena that the drive-will in every age,
in each sex and in every direction can put forth;
they are like meteors, they gleam; but the more they gleam,
the more readily will they burst. The individuals of the drive-will are with
all their force of will not free, but bound to nature: they are in the power of
dim powers. True, they do not rush on blindly, for if they
blindly followed the promptings of the passions, they would become eccentric,
all too comet-like, and then we lose interest in them. But
even when they wish to keep to the track of the Idea within certain
limits, their way is the most dangerous of all. They have
erectness, soaring, movement and speed; but the foundation
sways, and the movement is like a dance over the abyss.
They can gain power over others, but none among them
all can in the stricter sense gain power over himself.

It lies in the drive-will that it is moved by something
over which it does not itself rule.

As now in the erotic. Falling in love, it is said, is
something that comes upon people without their actually knowing
whence and why. It is not only youth that is
exposed; no age is quite safe: half-old men especially,
when it comes upon them, are as a rule badly shaken.
But what comes thus of itself can also pass
away again of itself; it is the whims of the drive-will, no one
can help it. The uncertain shows itself especially when two
conflicting feelings or two passions tower against each other,
like two waves, and a surf arises.
The person concerned is as if in distress at sea; he does not command
which wave shall become the highest or, in other
words, which of the two wills in him shall become the
strongest. If the erotic will is set, for example, against the artist's will,
then one has a beautiful tale of a chivalrous
% --- p. 20 ---
\setcounter{footnote}{0}%
\opage{20}lover who burned his harp to save his beloved from
dying in the cold. But the relation can easily be reversed; for
many a poetic lover has let his beloved freeze to death, and
preserved the harp. No one is responsible, the faithless one as
little as the faithful; no one can be called to account, since
everything depends, after all, on the stronger wave.
There is a Romance by Schak Staffeld, ``The Return''; it deals
with a Moorish warrior who returns home from the banks of the Tagus
with ``the dust of battle in his locks and its blood at his
steps;'' he calls to his beloved: Selindaja! But she
makes evasions; he has to name what he wants, the grotto,
the almond tree, the lattice; there is constantly something in the way, everything
has become different. At last he bursts out in noble wrath:
``Ha, faithless one, while I gathered these fifteen years of honor, without
shame thou hast given faith and heart to another;'' but she
answers him coldly, like one who has a good conscience:
``My heart I have not given away, Allah hear my words! Almoralif
has taken it, and what could I do about it?''
% DA: Kampens Støv i sine Lokker og dens Blod ved sine Fjed;
% DA: Ha, Troløse, mens jeg sanked disse femten Hædersar, uden Blu Du til en Anden Tro og Hjerte givet har;
% DA: Hjertet har jeg ei bortgivet, Allah høre mine Ord! Almoralif har det taget, og hvad kunde jeg derfor?

That the heroes with all their strength of will are at bottom
just as uncertain of themselves as the lovers, --- provided that
their will is only a drive-will --- can easily be proved. If one examines
what it is that actually holds the whole together in a hero-character
with so many plans, under so many oscillations and
crossings of external circumstances, and establishes its unity,
then one will find on closer inspection
that it is by no means a self-potent will, but some
pervading, ruling passion. Such a
passion does not always appear on the roaring
surface, it prefers to keep beneath the waves and works
quietly in the deep. But that it is not the hero himself, but the passion,
that is the actual master becomes conspicuous
when circumstances so contrive that a
new and mighty passion is awakened in the mind, so that the deep is stirred and the
two adversaries come into conflict with each other. In such
a case the strong man can suddenly become weak, get
fits of madness, and the character, as it were, go to pieces.
From ancient history just a single example: Marcus
% --- p. 21 ---
\setcounter{footnote}{0}%
\opage{21}Antonius. In the triumvirate between Octavian, Antony and
Lepidus, --- for vengeance on Caesar's murderers, Antony could,
as regards endowment and talent as a general, certainly in all
respects measure himself against Octavian. Nevertheless he succumbed
when it later came to enmity between them.
Why? The common explanation is that Antony gave
himself up to dissolute passions, while Octavian
in the confusion of the struggle understood the art of governing
himself, of calculating each step on the path of propriety, and of keeping
the folds of the toga closed. One would, however, be very much
mistaken to suppose that Octavian was a man who only
followed the commands of reason, while Antony let himself be led by
blind impulses. The truth is that both stand under the
drive-will. The difference of the characters is essentially to be determined
by the relation among their passions. Octavian, the
later so fortunate Emperor Augustus, had the good fortune of being
governed by only one passion, a passion which in
him was strong enough to compel and cow the rest. Octavian's
health was that he had only one sickness --- the lust for rule.
Antony, on the contrary, had the misfortune of being seized
by two passions, infected with a double sickness, lust for rule
and lovesickness. He came to Egypt, saw Cleopatra, and
--- what could he do about it? --- she conquered him, made
him her prey. Then the battle of Actium became decisive.
Antony fled because Cleopatra fled --- to Egypt,
and when she, the faithless one, kept herself hidden and had it spread
abroad that she was dead, Antony took his own life.
The hero fell, an unhappy victim of a love as faithful as
it was irresponsible.

By its nature the drive-will is the same, whether it
expresses itself in the small relations of everyday life or manifests itself in
great world events; but when we wish to
see through its essence, it is best, for the sake of clarity,
to see it in broad outline. It then appears that
the secret of strength is weakness, that independence is
dependent, that effervescent freedom succumbs to
necessity. The heroes and heroines of the drive-will go
% --- p. 22 ---
\setcounter{footnote}{0}%
\opage{22}boldly on board the stream of events; it makes a
splendid show to see them under sail: every vessel goes under
full sail, for the passions fill the sails. But here are
dangerous shoals, the compass has incalculable deviation, the drive-will
is like a drunken helmsman, and the sheet anchor is missing.
So they suffer shipwreck, they run aground on the rocks of fate.
Against the restless element of the drives, the inclinations, the passions
stand the barriers of existence, the strict laws of the world,
the firm walls of reality; it is the power of fate, ---
so it ends tragically. That a single individual, fortunate like Octavian,
escapes shipwreck is only a deception, a phenomenon
that concerns only the external; in the matter itself they are
all under fate.

It is the clash between the restless and the firm,
between the drive-like and the rational, in which the self is dissolved
and perishes. When the external resistance incites
the passion, when the inner seething and boiling increases, and
the surf rises, then the deep has its own thoughts, it
thinks of destruction; the great ships, the proud sailers
are dashed to wrecks, the forces of the drive-will are scattered. The drive-will is
thus not the free will; if the will is to become free, it must be
transformed, raised to a new stage and subjected to the laws.
In what follows we shall see an attempt.

% --- p. 23 ---
\setcounter{footnote}{0}%
\lecture{Third Lecture.}\opage{23}

What takes place with the will, as it raises itself from its first, natural stage up to its next and higher one, and from
drive-will becomes \emph{reason-will}, is a transformation. That this transition is more than a change, that, essentially
viewed, it rests on a transformation, must be expressly emphasized.

The drive-will can, as has been shown in what precedes, be cleared, altered, idealized; from being a sensuously blind will
it can become a seeing will, a will with thoughts, ideas, plans; that is an advance toward the better, an ennobling, a great
change, but it is as yet no transformation. The innermost mover is drive and remains drive; the whole falls within the
domain of unfree, obscure natural impulses. So long as the drive rules, passion is the proper measure of strength; there are
not only flaring, hot-tempered, violently erupting passions, there are also steady, silent passions, brooding, dreaming,
smoldering passions. But within the whole circle of immediate natural impulses, whether they be lower or higher, hot-tempered
or steady, this holds: the greater the passion, the stronger the will; the strength of passion is the condition for the force
of the will.

Only as the will passes over from being a rational drive-will to becoming what we must call, in the proper
% --- p. 24 ---
\setcounter{footnote}{0}%
\opage{24}sense, reason-will, does the transformation take place. Everything becomes different from the ground up. The drives
are not eradicated, but their power is taken from them, they are tamed; the natural impulses enter into the service of
reason, and must pledge themselves to unconditional obedience. New motives, a whole crew of new motives, come to the helm;
it is,
literally speaking, a new dynasty that ascends the throne. In the place of drive steps insight; not obscure stirrings, but
clear thoughts, not the enticing images of a conjuring hope, but cognition of the universally valid, firm conviction of what
is true and right in and for itself, is what from now on is to move the will. Here no one asks about pleasure and
displeasure, not about inclinations, but about duty and right, for duty and right are personal expressions of irrefutable
laws. Here in this view there is a loftiness that impresses; a cool sanctuary, a dome with eternal stars arches over what is
temporal. Man grasps his essential vocation and gains respect for himself. But in what attracts there is at the same time
something repelling, something that can awaken a shudder. To make the will free, reason must begin by cutting the bonds that
bind it to the maternal root of the drives; every nerve-fiber, every strand and thread from the mother stem must be torn
asunder. If reason is to rule, it must rule in One and All. Such a self-ruler does not come to terms with the stirrings of the
heart, will not parley with capricious inclinations, will not recognize any veto of feeling: reason is a sovereign with
absolute monarchy.

It will be understood from this that what matters essentially is not the externals of the action but its internals. An action
that springs from a rational drive-will can in its outward aspect bear a striking likeness to an action that has sprung from
pure reason-will: the difference lies in the inner mover. It is, when one speaks of historical persons, often difficult to
pronounce a sure judgment. An Agamemnon sacrifices his daughter, a Junius Brutus ``slaughters his sons on the altar of civic
% --- p. 25 ---
\setcounter{footnote}{0}%
\opage{25}freedom'': what sort of actions are these, really? Those who are at once at hand with their: ``blind fanatics,
inhuman monsters, unnatural fathers!'' betray precisely no deep insight either into the nature of drive-will or into that of
reason-will; for it is certain that there are scenes, collisions, terrible clashes of conflicting demands, in which cold
reason can tyrannize over the feelings of the natural heart fully as strongly as a drive-will glowing with passion.

The actions cited can, when seen at a distance and considered in and for themselves, be actions of reason, for it is indeed
the command of reason that everything individual shall be sacrificed to the universal. Yes, they \emph{can} be that. But if
we look more closely, then neither Agamemnon nor Brutus stands on the standpoint of pure reason; they are both absorbed in
finite ends, even if aimed at the universal of society: their will is a potentiated drive-will. In the domain of finite ends
reason is indeed to regulate the drive, but for the rest to spare it, so that within certain limits it is allowed to be
itself in all naturalness. From this point of view a dark shadow falls over the great figures, Agamemnon and Brutus; for the
father's heart has just as full its indisputable right as the universal of society has its claim upon the king and the
statesman. In the name of drive-will we demand, and with justice, that reason shall be human, and that its stern commands
shall let themselves be reconciled with humanity. The Greeks demanded that too --- the Spartans excepted, however --- but
not the Romans.


The remarkable difference that exists, on the whole, between the Greek and the Roman mind, between the Greeks' vitally warm,
brilliant drive-will and the cold but consistent Roman will, is recognizable even in the examples chosen at random. The
harshness, the revoltingness in Agamemnon's act of sacrifice the Greek legend has in its own manner known how to soften and
reconcile with beautiful humanity. Artemis, the angered goddess, comes
% --- p. 26 ---
\setcounter{footnote}{0}%
\opage{26}to the aid of human feeling; she demands her victim, but she also saves it; she saves Iphigenia, carries her away to Tauris and appoints her her
priestess. The impression of the divine mildness removes what is revolting and works like a bath of ideality in which the
natural feeling is purified. Otherwise with the Roman scene. To be sure, Brutus's sons had offended against the Republic,
inasmuch as they belonged to the young conspirators who wanted to reintroduce the monarchy, and Brutus, ``the liberator of
the fatherland,'' could not without self-contradiction vote for their pardon. But the manner in which the execution is
depicted, the petrifying coldness with which the father, without even turning his face away, sees his sons' heads fall, is
revolting to all humanity, a symbol of the Roman will's contempt for natural, human feeling.

This heartless harshness, this inflexible consistency that stifles all natural feeling, can easily impress the beholder,
confuse his judgment and mislead him to regard the Roman will not merely as more rational than the Greek will, but even as a
pure and great and exalted reason-will, an exemplary, systematically carried-through statesman's will. That is an illusion.
It is quite true that the Roman will is systematically consistent in a different way from the Greek will, and therefore also
stronger. The Greek will in its spring-fresh youthfulness has forces, great, brilliant natural forces; but the Roman will is
more than a radiation of forces, it is a peculiar unity of forces, a power, a world-conquering, nation-devouring will to
power. That this will to power, far from being what we must in the proper sense call reason-will, is instead a spiritless-spiritual
natural will, a wolf-greedy, insatiable drive-will, which, when it has devoured everything, must end by devouring itself,
needs no proof in this context: history furnishes the proof.

The modern concept of reason and reason-will is as yet no clearer, in a practical sense hardly as clear, as it was
among the sages of antiquity. We demand reason and
% --- p. 27 ---
\setcounter{footnote}{0}%
\opage{27}we demand humanity; but the question is whether these two demands can really be united, given that reason, mind you,
is to be our One and All, the sole, highest governing power in our humane world.

We demand that reason shall govern us and yet spare the inclinations of our heart; we think that this can well be done,
since our noble inclinations after all bear witness to rational predispositions and therefore must be willing to enter under
the unconditional dominion of reason-will: that may nevertheless be a mistake. We demand of reason that it shall not only be
strict but also fair, shall take account of our frailties, in short: bring us up as a strict but loving father brings up his
children. But the simile limps. Reason is no loving father, but an inflexible ruler, cold, unchangeable like the eternal
laws. When the father chastises his little boy, the little one repents, promises penance and amendment: ``will be good, will
be sweet!'' Then the storm passes over and it comes to a reconciliation, although the father knows that the remorseful
sinner can as easily forget the promise as he forgets the stinging rod. But thus reason cannot treat % print: bebandle (misprint)
us who are of age, cannot possibly at one and the same time both scold and forgive us rational beings. We demand that reason
shall respect our individuality, shall show accommodation, be a pliable, humane and tractable reason: on these constitutional
conditions it shall be allowed to govern. But on these conditions reason will not enter. And why not? Because it cannot do so
without contradicting itself and becoming unreasonable.

Is reason then despotic, inhuman and cruel? By no means. Far be it from us to wish to lower reason, to speak ill of it and to
try to put a stain on its pure, exalted character. We will, on the contrary, let reason itself speak. It says to man: You have
to choose. Will you remain standing at the standpoint of drive-will, well and good! I shall serve you with all that I can,
with counsel and admonition, instruction and
% --- p. 28 ---
\setcounter{footnote}{0}%
\opage{28}enlightenment in every respect. But then you must know that I am only your servant, not your master. You may follow my
counsel, if it so pleases you, and you may reject it if you will: I do not answer for the consequences. If your feelings
become discordant, if your passions rise in revolt, if your mind comes to a boil and a maelstrom forms that sucks you into
itself and grinds you down into an abyss, that will be your own affair: I am without responsibility. If, on the other hand,
you wish to be delivered from your nature-bound condition and lifted up to the stage of independence, then we must reverse the
relation; then it is I who am to command, and you who are to obey, then I am not your servant, but unconditionally your
master: choose now! That one cannot call inhumane; it is, as it seems to me, quite rationally spoken by reason.

The only thing one can fear is that reason, when it becomes so autocratic, will forget to make use of the fair royal right
that is called the right of pardon. And so it is. Reason is among all absolute rulers the only one that knows nothing of a
right of pardon. Reason is the selfless universal whose essence and content is exhausted in the coherence of laws. Reason-will
is the self's own will, inasmuch as the thinking self alone wills that the laws of necessity be fulfilled. Liberation consists
in the self's willing what it must will out of conviction: dissolving itself into the universal and selfless.

At the stage of drive-will the individual is well off when his individual claims, be they material or ideal, are satisfied;
therefore health is a good, sickness an evil, wealth a good, poverty an evil, etc. With that sort of satisfaction reason-will
cannot concern itself. All goods, both inward and outward, both non-sensuous and sensuous, inasmuch as they merely serve to
sate the immediate individuality, are of subordinate significance. One need not precisely despise them; but they are subject
to chance, are conditioned by fortune, and must on that account be reckoned among what is in itself indifferent. The sage is
raised above the vicissitudes of fortune; he hopes for nothing,
% --- p. 29 ---
\setcounter{footnote}{0}%
\opage{29}for he fears nothing: that is the old Stoic ideal, and this ideal of reason is as yet the only consistent one. The
sage is sufficient unto himself; he wills what Fate wills, his life is a revelation of Fate: so he is strong as Fate,
rock-solid, unshakable as Fate.

If we now consider this ideal of reason --- and there is no other --- somewhat more closely, it turns out that the
much-praised agreement of character with itself, of which so much has been made in ancient and modern times, dissolves into
contradiction upon contradiction, and that these contradictions by no means lie in the eternal reason, but in men's own
misunderstanding of what reason is and for what it is to be used.

It is then, in the first place, easy to see that the ideal set up cannot be realized. The high-flown phrases, that the sage
is sufficient unto himself, is happy even in a glowing furnace, is king, priest, seer and poet, must be taken for what they
are --- phrases.
If it is said:
that such an ideal is set up as a goal toward which one strives in the infinite; it is reached only in infinite approximation:
then it is not difficult to convince oneself that the approximation contains a contradiction. For so long, namely, as
everything that the natural individuality must regard as essential goods: health, talent, wealth, reputation,
power, influence, etc., has not yet been reduced to something altogether indifferent, and so long, therefore, as the inner
mover still remains an immediate drive,
so long man continues to stand under the drive-will.
On this stage there can be as little talk of an approximation to the pure ideal of reason as, by traveling around the earth,
there can be talk of an approximation to the moon.
If one makes it a serious matter that the conscious reason shall in all actions be the inner mover, so that no one dares
desire anything unless reason commands it, no one hunger or thirst, love or hate, except by the express command of reason,
then the inhuman begins.
% --- p. 30 ---
\setcounter{footnote}{0}%
\opage{30}But a goal that is in its innermost essence inhuman cannot be an ideal for human beings.

The contradiction appears in all real actions. That the temporal goods without exception are indifferent means, in other
words, that all finite ends are indifferent. The sage must be able to find his way in all positions, to be active under all
circumstances. Reason-will must act, its vigorous essence is acting. The sage can, e.g., be a general. Then he fights a
battle. Why? To win it! Yes, so a fool thinks, an individual who still finds himself under the drive-will. No, the sage
fights the battle for the sake of reason, solely and alone for the sake of reason. Whether he wins or loses, comes away from
it with honor or with shame, is indifferent to him, wholly and utterly indifferent. But when all external aims are
indifferent, then the corresponding actions must surely be indifferent too. Why then should the sage condescend to act?

The sage works and acts, not for the sake of other human beings, not for the people's, not for the state's sake, but solely
for his own sake, that in his being and working he may know himself to be in agreement with himself and with reason. From
this it follows that his measure of virtue and vice must become a different one from that which is adjusted for civil
society. For since every external action is in itself indifferent, the character of the action, its worth or worthlessness,
must rest exclusively on the agreement of character with itself in the agent. The sage can by virtue of his demand for
self-agreement be entitled to commit actions that in the eyes of men look like crimes; he can rob, murder, steal, lie, all
of it in consequence of pure reason. In such respects nothing is forbidden, except the impossible.

But, one objects, if the pure man of reason can become such a monster, how can characters as noble and exalted as a Socrates,
an Epictetus, a Marcus Aurelius, be counted among men of reason?
% --- p. 31 ---
\setcounter{footnote}{0}%
\opage{31}We answer: They stand on the stage of reason-will, inasmuch as the inner determinant is not an immediate and obscure
drive, but a truth cleared in thought. But by an ethical self-deception they have made reason into something more than it is,
into a spiritual self-being, a willing deity; they have out of their own rich interior put something into reason that
cannot possibly be derived from its essence. They are characters with a dawning, religious background: the personal
presupposition is: that the self is eternal. Such a presupposition is indeed not in conflict with reason, but
--- it must be impressed again and again --- it cannot be grounded by reason, cannot be derived from the pure concept of
reason. The concept of reason, held fast by itself, is the concept of the selfless universal: grounds of reason are grounds
of necessity, not grounds of freedom.

The determination of freedom that can be derived from pure reason is as self-contradictory as possible. The sage has a will,
for he has a self. So he begins then by purifying this self, pouring it into the crucible of reason in order to free it from
dross and impure admixtures. He gathers himself and takes hold with a strength of will, an energy, as if this self were
destined for a whole eternity. But precisely here the contradiction reaches its height; it culminates, so to speak, in the
zenith of reason. For the self, seen in this light, is not eternal; it is perishable like the things in nature. The existing
self is individual, and the individual dissolves and perishes: only reason is eternal, eternal as the laws. When the sage has
played his part on the world's stage, no matter whether it is the beggar's, the tailor's or the emperor's part, and sees that
the part is over, then he steps off. He need not wait for natural relief, he can relieve himself, he can freely extinguish
his light and go away: to give himself death is the expression of the highest will of fate, the determination of will in
which the self itself dissolves itself, the last act of reason of formal freedom.

But when man's self is perishable like everything in
% --- p. 32 ---
\setcounter{footnote}{0}%
\opage{32}the world, and what is perishable is an indifferent, then it is surely evident from this that the self too is an
indifferent, its independence and agreement with itself at bottom an indifferent: and yet it was precisely for the sake of
the self and independence --- that is, for an indifferent! --- that one was to live, work and act in the last instance.
This is the consistency of the pure reason-will carried through in One and All.

Modern consciousness admits that the genuine Stoics entangled themselves in contradiction, but must, naturally!
grant them right that the ethical ideal of reason is the only true one. The matter is that those old thinkers were
one-sided. Their view of the world was cold and barren: the content of existence, the luxuriant growth of the tree of life,
its branching trunk, the playing light, the lush shoots of individuality with leaf and flower in fragrance and color, all
this is missed; all this was overlooked, for the old thinkers were, as said, one-sided. But the new thinkers are all-sided,
they take everything in, the smallest with the greatest, and overlook nothing. How do they go about it? It seems
complicated and yet goes straight: let one simply make nature into reason, and the content of life --- the concrete --- comes
of itself. Reason is natural in feelings, drives and passions; nature is rational in thoughts, ideas and laws. The demand of
reason is the right of nature: reason in flesh and blood, nature's reinstatement in its right, is the true ideal of the
future, the deeper meaning in ``the Gospel of the Flesh.''

So far the ancients could not reach, their simplicity was still too great for that. To be sure, they too began with a general
unity of reason and nature: to live in agreement with reason was with them to live in agree-.
ment with nature. But this held only quite in general; as soon as they came to the particular, the split set in. To make the
temporal eternal and the eternal temporal, the indifferent essential and the essential indifferent, is a task whose solution
surpassed the old
% --- p. 33 ---
\setcounter{footnote}{0}%
\opage{33}thinkers' powers, and must on that account be reserved for the new thinkers.

The dominion of reason and the Gospel of the Flesh! Dare we now also give ourselves over to the bright hope that so colossal
a contradiction will really be resolved? Instead of going forward, we go backward; instead of raising ourselves up to the
higher stage of reason-will, we sink imperceptibly down below the drive-will. Reason is not modern, any more than nature
is; it is not at all said that modern consciousness and eternal reason are One and the Same.

The stage of reason-will is a necessary transitional stage, and cognition of practical reason an essential condition for a
strong will. What purifies and educates lies in this, that the human self learns to bow before something universally valid,
learns to resign itself and through infinite resignation receives an impression of something eternal. But reason cannot
teach man to grasp his own individual self as eternal. If it could, then there would have to be given in science a fully
valid, irrefutable proof of the immortality of the soul. But that such a proof has neither been given nor will ever be able
to be given is a clear matter, once we have first learned fully to see what reason is, and why its eternity must be selfless:
the possibility of immortality can be proved, but its reality cannot be proved.

If one wishes to prove the immortality of the soul on moral grounds and appeals to the claim that a complete development of
our free moral personality would be impossible if our soul were not immortal: then one oversteps the boundaries of reason,
overstrains its powers and makes a personal stake for whose general validity pure reason can neither will nor do go surety.
So it remains the same as before: reason, the universal essence with ideas and laws, is eternal; man's individuality with
its organically conditioned self-consciousness, its drives and requirements of life, is temporal. If the decisive weight is
laid on the latter,
% --- p. 34 ---
\setcounter{footnote}{0}%
\opage{34}then the drive-will rules; if self-determination proceeds from the former, then the reason-will rules.

Drive-will is a stage for itself and reason-will a stage for itself. The two stages cannot be screwed together into one. One
cannot take pieces out of the drive-will in order to mend the reason-will, any more than the reverse. Personal independence
cannot be consolidated by gluing something individual that is not eternal together with something eternal that is not
individual. The freedom of the will is no patchwork: one can patch an old shoe, but one cannot patch ideals together. If
reason-will is to be united with drive-will, then they must both be transformed and taken up into a third, higher: this
third, higher is spirit-will.

% --- p. 35 ---
\setcounter{footnote}{0}%
\lecture{Fourth Lecture.}\opage{35}

We live in the age of discoveries; here there is constantly asking and hearing of what is new: new machines, new stars, new
solar systems, new sulfur matches! who can count it all; nothing great or small escapes the investigator's gaze. It swarms
forth, the hitherto unknown, from the depths of the sea, from the depths of the firmament and from the bowels of the earth.
What is scattered is gathered, ordered, compared, interpreted, evaluated, examined. It is the content of reality, the
coherence of the world, the forces of things and the realm of laws that the thinking, observing reason tirelessly explores.

The great voyages of discovery go not only outward but also inward. Reason does not dwell on the merely sensuous; it asks
about the non-sensuous, it wants to discover the eternal in the temporal, to recognize in all that is becoming that which
is. What, then, is this eternally existing? Is it the basic parts of matter, unchangeable atoms, of which all things are
composed in order to be dissolved again? Or is it a superconscious and precisely therefore unconscious being of thought,
nature's blind, indwelling wisdom, which works and weaves in all phenomena? Or is the secret perhaps an infinite Nothing,
the Nothing that levels All, the dark Night that hides All, the holy silence in which all voices fall mute, the overfull-empty
abyss, the eternal Ginnungagap, from which All proceeds, and to which All, the Whole, shall at last return? What is nature,
what is spirit, where is the bond that binds them together? Here reason makes discovery upon discovery; the one worldview
chases the other, the
% --- p. 36 ---
\setcounter{footnote}{0}%
\opage{36}one philosophy displaces the other, the one thinker outbids the other. Materialism, spiritualism, realism, idealism,
pantheism, atheism, nihilism: it is a motley, sprawling sample card that lies before us. And what does this prove? That here
are many, and, when we look more closely, coherent attempts, all of which strive toward One goal. Reason wants to clarify its
view on all sides; free thought will not let itself be bound to any dogma, will not yield to any dictum of authority, will not
turn aside for any arbitrarily set up barrier. --- That is right, so it ought to be. We should not set the light of reason
under the lampshade of traditions, not put free thought in chains, not wish that the many and great discoveries might at last
come to an end. On the contrary! The more discoveries of reason are made, whether they belong to the outward or the inward
world, the clearer the knowing consciousness must become.

One discovery remains, however, a single one, but very essential, a discovery that is constantly overtaken in learning, and,
despite the much learning, hastily treated; its riddle lies in the question: what is a \emph{human being}?

``A human being,'' answers modern knowledge, ``is a happily organized animal; that is plain from the developmental history of
the globe. Man descends from some animal form in the primeval world, probably from an old pregnant, ambiguous ape.'' ---
If by this descent only a zoological-natural transitional link is to be thought of, then the inquiry concerns natural science
exclusively, and we may let the hypothesis stand for what it is --- a hypothesis. But if the descent is to be understood in
such a way that human nature is originally animal, an ape nature, which in the course of thousands of times millennia through
countless changes and gradual transitions has worked its way up to become a self-conscious rational being, a being whose
obscure, nature-bound ape-will has little by little become a humanly clear and free reason-will: then the discovery becomes
philosophical, and falls under the doctrine of spirit.

% --- p. 37 ---
\setcounter{footnote}{0}%
\opage{37}The philosophical, it is said, is so hard to understand; it is only for the good heads. Yes, the philosophical that we have
before us here is quite certainly hard to understand. An ape nature that is conscious of its reason, scientific apes, apes that
philosophize! the philosophical in such a philosophy lies beyond common human understanding; it can only be understood --- in
clear logic --- by the good heads of the family.

If one wished from this point of view to go into examining the value of the yield that man has won through the gradual
development, then there are here, seen from one side, certainly astonishing results to show; but if we ask how far all this
has really been to the benefit of the race's freedom and happiness, then misgivings arise, not least with regard to the will.
The animal wills what it must will, it is content with its drive-will; man has taken upon himself to will both what the drive
wills and what reason wills: that is a great inconvenience. The ape is happy, it always prefers what it ought; man must make
a distinction between what he prefers to will and what he ought to will; he is not happy at all. By his admirable
self-development, his much culture and civilization, man has brought upon himself the great misfortune of being burdened with
two wills that cannot be reconciled.

That is what comes out of the many discoveries. Man has discovered himself and in the self discovered the split; the split is
in his interior, the contradiction goes to marrow and bone, for the self is animal, and yet rational. The ape, which has never
engaged in discovering itself, has happily escaped the split. The animal has a natural self, but never sees it, and gets off
easily from all trouble with the contradiction. Man, who has turned his eye inward and seen his conscious self, has become a
victim of the contradiction. From the contradiction doubt springs, and with doubt follows error: man has discovered doubt and
in doubt the possibility of error. In older days a detailed register was kept of the many positive, both bodily and mental,
% --- p. 38 ---
\setcounter{footnote}{0}%
\opage{38}advantages that were to distinguish man above the animal. The positive in all these advantages recent criticism has
as good as erased; the negative advantages, on the other hand, criticism has not been able to erase: man can ruin his life,
distort his task, go astray and miss his goal; the animal has not that advantage. Man has discovered virtue, for he has
discovered vice, and he discovered vice precisely on the day he discovered duty. The animal is virtuous, but without knowing
it, it acts in accordance with its nature, it knows not vice, for it knows not duty. Man has made the acquaintance of duty,
alas! it makes a troublesome notch in virtue: if there were no duties, then we should all be virtuous. Man has discovered
death, a gruesome discovery. That the animal must die just as well as man is certain; but death as death does not exist for
the animal. Death is in itself as a nothing; but it is reserved for the self-conscious being to think this nothing, to meet
this nothing as a power and give it existence in thought. The animal lives in the moment and enjoys in the moment; every given
moment gets its own, and the next moment does not exist. Man sees in the present moment the possibility of a next moment, and
in this possibility he sees death. So the sight of a next moment can suddenly disturb the enjoyment, the admonishing shadow
can show itself in the midst of the feast, like Banquo's ghost at Macbeth's banquet. Man has discovered eternity; that is
still the very worst. This discovery sets the crown on all his advantages, yet, mind you, negatively, for it is a crown of
thorns. By becoming conscious of the eternal, man discovers that the self-contradiction in his being is insoluble: the
eternal obligates, and the contradiction is in the obligation.


To what is man obligated? To do the good --- for the sake of the good. In this lies that the good may not be measured with
the yardstick of finitude and ranked with the merely useful, the expedient of the moment; a bare and naked doctrine of
utility is in the moral sense an
% --- p. 39 ---
\setcounter{footnote}{0}%
\opage{39}impudent denial of the eternal. But when the good is to be willed because it is good in itself, then a human being can
surely in the temporal sense become unhappy, solely because he wills the good. When, then, our descent from the animal is
understood in such a way that man dies like an animal, and the meaning is that the human self is as little eternal as the
animal one: then the human being who wills the good for the sake of the good must surely entangle himself in an insoluble
self-contradiction.


The good, it is said, has its reward in itself; man, who feels the obligation, gets an elevated self-feeling, an involuntary
sensation that the eternal is a present. That is an ambiguous speech. It can be true, so long as the question about the
self's eternity still stands unanswered. It is a bad sign of a sense for the good when a human being wants to defend himself
against the claim of obligation by demanding proof that a temporal sacrifice will really pay off some time in a coming
eternity. But the speech becomes untrue when it is discovered at once, and both things are seen equally clearly in the
discovery: both that the self which is obligated is temporal, and that the obligating good is eternal.

As the living being is, so must its basic demand, its essential need also be. The animal lives in temporality, it demands the
temporal, and only the temporal has claims upon it: it has no contradiction of obligation to contend with. Man, upon whom
there rests a serious claim both from the temporal and the eternal, is placed between a double fire. By keeping to the
temporal he must surely, in the essential sense, be able to satisfy his innermost self, as truly as this innermost self is
temporal, only temporal. But then comes the thought of eternity with its unconditional good and demands that the temporal
self shall in this unconditional good recognize its true essence, shall, although it knows within itself that it is only
temporal, obligate itself to live as if it were eternal; this contradiction cannot be resolved. The being that obligates,
% --- p. 40 ---
\setcounter{footnote}{0}%
\opage{40}and the being that is obligated, must in their original ground be of the same kind, otherwise the obligation is a
matter of compulsion without moral worth: only an eternal self-being can be obligated to the eternal. To demand of a temporal
self that it shall let itself be obligated by an eternal good is just as contradictory as if one wanted to force an animal
that is organized only to move in the water to fly in the air.

Unless the many discoveries are to end by making existence meaningless and life intolerable, man must once more see to
discover himself: he must discover faith and in faith discover spirit.

But discover faith! Here modern consciousness meets with its philosophical. To discover faith, it is said, is even more
superfluous than to discover gunpowder; faith was discovered long ago, and its secret seen through. Faith is no new star, but
an old horn lantern, by whose magic glow anyone can see what his imagination would most like to beguile him with. No, not
faith, but knowledge! Knowledge is the new star, knowledge is the sun of the world, the royal light of day, that beams in the
eye of the age. So runs the announcement. What do we get to see by the help of this light, when we, mind you, look into the
self? Doubt and doubt and doubt again. The self in doubt and the self in contradiction, eternity in thought and death in the
heart: consuming self-contradiction! Then the light of knowledge at this point cannot possibly be the true light. Here, moreover,
there is not one light, but many. The waves of light roll so wildly that the many gleams partly intensify and partly darken one
another. They do not come from the heaven of spirit, but from the hiding-places of contradictions, from restlessly shifting
thoughts. Shifting thoughts give only fluttering light: shooting stars and will-o'-the-wisps, weather-lights that flicker at
the gusts of doubts, such bewildering lights cannot possibly lead us to discover faith.

But has faith then not long since been discovered? Yes, to be sure!
But the matter is that with discoveries in the personal
interior, in the obscurities of the mind, in the self's hidden
% --- p. 41 ---
\setcounter{footnote}{0}%
\opage{41}depths, is otherwise than with other discoveries. When one man in science has made a discovery, others may well step in to
test it and, if it stands its test, appropriate it. But those who appropriate it in this way cannot give themselves out as
having made the discovery at first hand. Nor is that needed: what is discovered stands firm as a given, and what is thus
given can at second, third and fourth hand come to general knowledge, and its cognition spread in ever wider circles. But so
it is not with the discovery of faith and the discovery of spirit conditioned by it. Both discoveries, presupposing one
another, are subordinate determinations of one and the same fundamental discovery, whose secret lies in the self, and which
for that reason every human being must himself see to bring to light. The discovery of faith is a personal task for each
single individual, always a discovery at first hand, a primitive self-discovery. That this discovery must have its
presuppositions in the revealed religion is certain; but for that reason it does not cease to be primitive. What does it
profit, pray, that religion with its holy facts is given to us, when we cannot cognize its truth. To cognize the truth of
religion is to discover religion; but to discover religion is in turn to discover faith.

How superficial and loose the common conception of faith is, makes itself involuntarily known in corresponding phrases. If
it is external things, historical events, of which we do not yet have full certainty, then it is said: I do not know it, but
I believe it. If it is a scientific doctrine over which there is dispute, then it is said: I should like to believe that what
is set forth is right, but I do have a little doubt; in order that I may know that the proposition is true, I must have the
proof carried through. So it is with all those truths in whose nature it lies that they must be capable of being
demonstrated as universally valid, be it by immediate sense-perception or by inferences and compelling proofs. In all
scientific cognition faith counts
% --- p. 42 ---
\setcounter{footnote}{0}%
\opage{42}only as a provisional assumption beset with doubt; universally valid truths of reason will \emph{not be believed}, but
\emph{known}.

When now an empty, superficial mind or an individual taken up and captivated by science and scientific research wants the
same measuring-rod applied also to the religious, this is a striking proof that the knower has not yet discovered faith; for
in the religious sense faith is itself the highest human certainty. Faith here steps into the place of knowledge and for that
reason will not be grounded on knowledge. To know and to believe is in the religious sense the same: I know it, for I
\emph{believe} it. I know, for I believe, that this human self, which begins in time, has an eternal origin and bears in
itself its germ of eternity. I know that this self with its spiritual-natural individuality is created, not in the likeness
of the ape, but in the image of God; I know it in faith, as surely as I know by the senses that fire burns, that snow cools,
and the sun's rays give light and warmth. I know it because I feel it, and I feel it because I think it. I know it, not
merely because I feel it and think it, but because in all feelings, all conceptions and all thoughts I will it; I will it,
not from arbitrary whim or blind impulse; I will it because in calm, clear, deliberate self-cognition I must will it, and I
must will it because I love it. When this self-understanding is won, then faith is discovered, then the old saying holds: tell
me what you love, and I shall tell you what you are, for what you love, that you are, and what you essentially are, that
you love.

But --- it is objected from another side --- when faith yields complete certainty for all believers, and the doctrine of
faith can reasonably aim only at having convincing force for believers: is it not then evident from this that the content of
faith must be able to furnish material for scientific grounding just as fully as either the content of nature or that of
history; that, in other words, the investigator of faith must stand in the same relation to the facts of religion as the
investigator of nature to those of nature? --- It is an
% --- p. 43 ---
\setcounter{footnote}{0}%
\opage{43}old and stubborn, but for that no less mind-enfeebling error that we have before us here, an error that is quite otherwise
dangerous to life for the cause of religion than either materialism or atheism or nihilism or whatever else one may name.
There must indeed be a circle of cognition in which the truths of religion in their inner coherence and grounding come to full
illumination, a bright and clear circle of cognition in which the eternal laws of thought, with a rich apparatus of thoroughly
researched facts, continue to hold. But at the same moment one wants to place the doctrinal edifice of faith on an equal
foundation with what in the proper sense belongs to science, the whole edifice collapses.

What is to become of it when the investigator of faith wants to relate himself to revelation as the investigator of nature to
nature? Nature is an ever present, accessible to sense-perception and sensuous certainty. But where is there sensuous
certainty for faith? What is to become of it when the investigator of faith absolutely wants to relate himself to the
revealed facts as the investigator of history to history?

If the articles of faith are first to be grounded on historical testimony and accepted on external authority before they can
be believed, then Catholicism --- the papal, mind you, with the dogma of the Pope's infallibility --- is the only right one;
for outside the Catholic church it is surely just as impossible to give full historical warrant for three hundred dogmas of
Scripture as for three apostolic articles of faith. If, on the other hand, personal assurance of the truth of the articles of
faith is to be sought in living faith in the self-present Word of life, then this faith in the Word of life is a source of
cognition for itself, a principle that grounds itself in the self, independently of knowledge and science. Everywhere that the
disguised hybrid being of faith and knowledge still carries on its old haunting, faith is not discovered.

To discover faith is to discover spirit and in spirit to discover the will; for spirit is the ground of the will, indeed the will
% --- p. 44 ---
\setcounter{footnote}{0}%
\opage{44}itself. For the freedom of the human will in the power of spirit faith is the first condition.

The difference between spirit-will and reason-will can already here become recognizable. Reason-will is powerless against
nature, for it is only man who wills in the name of reason, reason itself is no willing being. Those who attribute to this
being a free, independent will must either set up reason as an idol or, on account of a confusion of language, mistake reason
for spirit. The divine Spirit has nature in its power; nature is created, borne and explained by Spirit. If anyone, on the
other hand, will maintain that it is reason that has nature in its power, then we can at once reverse the proposition and say:
it is nature that has reason in its power. Both views are equally true and equally false. It goes with reason as with laws:
such laws, such things; and conversely: such things, such laws. The essence of things expresses itself in the laws, but the
reality of the laws rests on the things; if we take the laws away, the things vanish, but if we, in thought, take all things
away, what becomes then of the laws? The mutual relation of dependence of things and laws, nature and reason, points to the
third, higher, that has both in its power, to Spirit.

Spirit has in the unity of its almighty will all the riches of nature and all the laws of reason; it does not move us as an
obscure drive, it does not rule as the cold law; it does not compel, it creates, it does not piece together, it inspirits.
It is Spirit that explains the individual in the universal, drive in reason, and reason in drive. Spirit wills the universal,
but it wills just as originally the individual. It does not dissolve the self, it does not efface the form; it gives the
eternal life and life eternity; it gives the living a self and the self
 form, a form consecrated to transfiguration: it fulfills the laws and it overcomes death.

At the stage of the life of spirit the will becomes strong; let us see what it is that gives it strength.

% --- p. 45 ---
\setcounter{footnote}{0}%
\opage{45}So long as the human will is referred to itself alone, so long as it, so to speak, must conduct the struggle on its own
account against the barriers of nature and all the resistance with which necessity and its old ally, chance, are able to
overwhelm a limited individuality, things look desperate for freedom: drive and reason both stand under fate. If the
struggling will can, on the other hand, get relief from an absolute strength, a helper who is himself will, almighty will,
then it clears up on the horizon of freedom, and the prospects become bright. The almighty will does not want to consume the
powerless will, but to take it up into itself and give it strength: the deeper this cognition penetrates, the more vigorous
the will must become. To be sure, in temporality there continues to be a discord between drives and duties: the drive-will is
selfishly vehement, the reason-will empty and cold; but in spirit and faith --- at the stage of spirit-will --- the split is
essentially removed. The disparate demands, the temporal and the eternal, begin to transform themselves and become homogeneous.
Reason is spiritual, it can be inspirited and filled with personal content; drive is warm with life, it can be cooled and
purified in inspiriting, in order to be warmed again. For spirit is not cold; it is the wellspring of all life, in blessed
longing with deep yearnings, it has all the warmth of inwardness.

God is Spirit, and God is love: Spirit is love. It lies in the essence of love that it unites the opposite and reconciles the
conflicting. Reason, that is to say: our rationally thinking consciousness, can see the truth in this proposition, so far as it
is held fast in all generality, but when the knot is tightened, and what is decisive is laid into love's ground, its aim and
its will, reason becomes blind. If reason is to assume an eternal will, then it becomes a general, unconscious willing, that
coincides with the eternal necessity, with nature's indwelling drive to work and to form. A divine will, which with
self-conscious aims has breathed life into the dust, has united an immortal soul with a
% --- p. 46 ---
\setcounter{footnote}{0}%
\opage{46}mortal body, and strained the opposition of the corruptible and the incorruptible so strongly that this our body must die,
in order that the soul may in truth live; a will of omnipotence, that has done all this out of love and with the fatherly
intention that man under the pressure of necessity and the burden of temporality shall raise himself up to freedom and appropriate
an eternal life --- such a free, conscious and personal will of love does not exist for scientific reason. We must find it
highly natural that those who seek all higher certainty exclusively in knowledge and science continue to regard a doctrine of
freedom of this kind as arbitrary and ungrounded.


Reason can see that the goal of freedom would stand high and be a magnificent goal for a strong will's striving, \emph{if}
there were a world to come. Yes, \emph{if}! But how do we get beyond this \emph{If}? How does one go about forming a clear
concept or even merely a probable conception of the constitution of a supernatural world? Which science undertakes to furnish
the proof of its absolute necessity, without smuggling in presuppositions that do not belong to knowledge?

The will of eternity that reason can cognize is in its unchangeable essence neither loving nor unloving; it loves nothing, for
it hates nothing. Love and hate belong, like longing and aversion, sympathy and antipathy, under the oppositions of finitude,
and fall upon the affects of individuals.

By noting the limit of human reason we do reason no wrong. On the contrary. To demand of a being what exceeds its capacity is
unreasonable; it is unreasonable to demand of reason what is beyond its powers. Human reason has made many wonderful
discoveries in the world and will hopefully make still more; but there is one discovery that it has never made and in spite of
all probing, brooding and prying will never make, that is the discovery of the self's eternal life.
% --- p. 47 ---
\setcounter{footnote}{0}%
\opage{47}That is precisely the discovery that lies for faith, the only discovery that faith in spirit and truth can make. Faith
discovers the self's eternity: it discovers spirit and in spirit love and in the spirit of love a blessed life. This
discovery is an inner personal experience, it is like breathing and must constantly be repeated; if breathing stops, life
wanes; if the experience stops, faith wanes, and with faith spirit and with spirit the vigorous will.

We live in the age of discoveries; we live to discover
--- the will of spirit.

% --- p. 48 ---
\setcounter{footnote}{0}%
\lecture{Fifth Lecture.}\opage{48}

\emph{When the spirit is strong, the will is} strong;
this proposition may be conceded without leading us
to any deeper insight, for the question is: whence
does the spirit come, and how does it gain strength
in man?

It is not with the spirit as with the drive. Whence
comes the drive? We need not look for it, it is
near to us, in our nature. It does not \emph{come}, it
is; it expresses itself involuntarily as it is, because it is.

It is not with the spirit as with reason.
Whence comes reason? It does not come from afar;
we do not stare up into the blue and await its coming.
It does not \emph{come}, it \emph{is}, in our being. In the same
moment in which we reflect on what the rational is,
we discover that reason is not something absent, but something
essentially present to us.

Can the same now be said of the spirit? To this one must
answer both Yes and No. Yes! The spirit is present;
it is indeed divine, and the divine is an omnipresent.
And yet No! for the spirit is not, like the
drive, immediately present in nature; it is not
necessarily present to thought, as reason is. When
the rational is thought, it is so with necessity; but
because someone thinks something spiritual, it does not yet follow
% --- p. 49 ---
\setcounter{footnote}{0}%
\opage{49}that the thinker has the spirit itself in living presence.
The spirit can thus be at once both
present and yet absent, can be infinitely near
to us, although we must nevertheless await its coming. It
is like a riddle. The secret is: that the spirit
\textbf{is not merely \emph{being}, but also \emph{will}}.

If the emphasis now falls --- the emphasis decisive for religion and faith --- on
the spirit's being will, then its
coming is a revelation, and the revelation a miracle.
Thus the spirit of the Gospel came upon the first
witnesses of the truth; it came as of a sound from heaven, as of a rushing
mighty wind, which filled the house, it came with
tongues of fire and with the gift of languages, it came
with the tokens of eternity, with the signs of omnipotence, it came in a
great miracle. But this miracle is precisely to knowing
reason a sign of offense; it is the miracle that has deepened
the chasm of divorce between the past and the present. How
everything has changed since that first breakthrough!
The spirit, it is said, is still the same; but where is
the power, where are its wonders?

The civilized present does not believe in miracles; the
rational cannot even hear them spoken of without
feeling nausea. The world, it is said, is no longer arranged
for miracles; miracles are against its constitution;
its nervous system cannot tolerate them, neither the earth's nor
heaven's. We turn our gaze downward, we dig
in the earth and find there below a great library of stones,
which in plain characters tell of the life and conditions
of past worlds; but in no region of the earth, at no point
do we find traces of a Paradise. We arm the eye and look upward:
stars upon stars! but where are the angels, where is
the throne of the Eternal? Heaven has opened, to
us too, but otherwise than in those days. The heaven of the present
cannot open as the heaven of the past opened for him who
ascended into heaven; our clouds could not bear him, if
he should ever come again ``with power and great
glory''. The believers see it, and yet dare not see
% --- p. 50 ---
\setcounter{footnote}{0}%
\opage{50}it. Faith is a veiled nun; she bashfully casts down
her eyes, but stolen glances pass through the veil. Therefore
--- so the keen-sighted assure us --- the generation is weak;
it has become knowing, and yet wants to be believing; its
being is divided and therefore weakened; it lives half in the
old and half in the new, half in the sorceries of
supernatural imaginings, and yet more than half in the
interests of reality with all the present
facts of a natural world: so the generation pines away in half-faith and
doubt, the will is enfeebled and characters grow weak.

It is quite certain that a serious crisis must be
gone through before the more recent culture works its way
out of that twilight of the gods in which its irreconcilable conceptions
of time and eternity, like giants and Æsir, visibly
fight the last battle. The spirit of revelation with
its wonders and nature with its unchangeable laws have
in our present consciousness taken up such sharply opposed positions
that half-faith goes about with a bad conscience and waits
for a final explanation. Thoughts of faith and thoughts of knowledge in
confused mixture, fantasies of religion and impressions of
science move in strange visions, misleading, like
the mirages of the desert to the traveler. The way through
this difficult going, the only way that leads surely to the goal,
is marked by a question that runs through it all, the question
of the will. For this question has the
peculiarity that it leads the strife about everything external
back to a single, decisive something internal.

What is the innermost mover of all existence: nature's
dark drive or an eternal, free, omnipotent will? That is the cardinal point
on which the whole turns. If the
first is assumed, the miracle is an impossibility, and historical testimonies,
be they never so many and never so concordant,
cannot breathe reality into the impossible.
If the latter is assumed, the miracle is possible. And yet
historical testimonies cannot give us full certainty of any
single actual miracle. Historical testimonies
give us only a probability; but the actual miracle,
% --- p. 51 ---
\setcounter{footnote}{0}%
\opage{51}it lies in its nature, is always improbable. Knowledge keeps
to the historically probable, and rejects the improbable;
faith goes beyond the probable, and seizes
the possible. The impossible cannot be believed, but the improbable,
which is yet possible, can be believed: here is the difference
between faith and knowledge.

In the realm of profane history the researching
science has a sure guide in known laws and a
measure of probability normed by them; when an event
is so probable that no one can deny it without
at the same time denying all historical probability, then
the event is secured as a fact. Otherwise
with the facts of revelation in sacred history.
Here the probable cannot be normed according to known
laws, for the miracle itself is the most improbable of
all. If faith now reaches out beyond the merely probable in order, in
the name of the omnipotent will, under the form of reality
to seize the highest possibility, then it seems exposed
to overlooking the natural-necessary connection of things,
forgetting the laws and straying into the fantastic. Against
this error stands the Word of life, the divine self,
the revealer of the will to power, firm, immovably firm as the rock
in the sea. His will is power, the secret of all possibilities,
he is the miracle itself, surrounded by wonders and signs.
And yet his own will, humanly seen, is as
an impotence; he has, literally speaking, no self-will: his
will is to do the will of omnipotence, the will of the Father, who
is in the heavens. Thus he stands on the ground of freedom, and
yet under strict laws. He shows himself as he is,
--- the truth! Yes, but in obedience. Not as an arbitrary
ruler, but as an obedient servant, he lives in the truth,
strives for the truth, suffers for the truth, dies
for the truth: he, the serving truth, does not break
the laws, he fulfills the laws. He does not found a new
religion by means of a new doctrine with a good example,
no, he founds what he is, and he is what he
founds, religion itself, \emph{the spirit-will} in flesh and blood.

% --- p. 52 ---
\setcounter{footnote}{0}%
\opage{52}It goes glibly enough with declarations like this:
``we live in the age of laws, not of miracles''; but is then
the highest concern of the human race settled by this?
Let us translate the declaration into another language, in order to see
rightly what it contains. We live in the age of laws, that
is to say: the will of omnipotence has gone out, only the finite
wills remain. The banner of our future bears the new
inscription: \emph{the will of the day}. With this it agrees quite
well that all wills of the day should come forward, and the will of the majority
sit on the thrones of power; but how this will of the day
will go about attaining the lofty goal: to fulfill
the laws and lift the masses up onto the steps of freedom, is not yet
explained.

The question of the miracle coincides in the last instance
with the question of the \emph{originality of the will}.
This question does not let itself be dismissed by appealing
to the unchangeableness of the laws of nature. One might then
just as well point to the unchangeableness of the laws of thought.
If the unchangeable laws of thought cannot even be said to
impair the human will, when the talk is of
new ideas, new beginnings and correspondingly new actions,
how much less should the unchangeable laws of nature then be able at
any point to impair the unconditioned freedom
of the divine will. Unconditioned freedom does not consist
in the will's being without conditions, but in its
having the conditions in its power, and being raised above finite
limits. Reason is the form for the content of wisdom,
but the content itself is determined by the will. The spirit has
not reason outside itself, any more than it has wisdom
outside itself. The omnipotent will can do
what it wills, act and move as it wills, it never
comes outside the infinite connection of the infinite whole:
it can indeed never come outside itself. But precisely
in the infinity of the connection freedom has its inexhaustible
possibilities.

For the omnipotent will all things are possible; that
does not mean that the omnipotence of the will can make the laws nothing,
% --- p. 53 ---
\setcounter{footnote}{0}%
\opage{53}raise contradictions to laws, make nature unnature, the circle
square, etc., but it means that it in and by
and with the laws can bring forth what in the eyes of men is the most improbable of all,
can, without breach of the hidden connection
of things, suddenly break all the measures of probability.
The omnipotent will, which has things in its power,
has also the laws in its power; but to have the laws in its
power is: to have their fulfillment in its power. He who
has thought in his power never breaks the laws of thought;
he who has nature in his power never breaks
the laws of nature.

There is a fundamental proposition on which the faith of revelation
with its assurance of spirit and will must be at one,
openly, sincerely, completely at one with the consciousness of the present,
namely the proposition of the unchangeableness of the laws.
When the validity of this proposition, rightly understood, is fully
recognized, then the connection between past and present
is spiritually restored, and the possibility of a mutual understanding
between reason and religion established. On
the foundation of this understanding rests the separation of that
which can be known scientifically, and that which cannot be known,
but must fall under faith. What is the original
ground of the unchangeableness of the laws of nature? ``The nature of things,
the rationally necessary connection of the world''; that is the explanation of natural science.
``The almighty will of God''; that is the explanation of religion.
Here there shows itself on the map of cognition
a distinct boundary between faith and knowledge, a
line of demarcation, quite as reliable as the one the Pope
drew when Columbus had discovered the New World.

The doctrine of experience begins with the facts,
keeps to the laws and halts at necessity.
Philosophy of nature goes a step further. It can, in the fundamental relation
of nature and spirit, point out the conditioned necessity's
issuing from freedom, the origin of the laws from the legislating
will. But with all this it must nevertheless remain
standing at the side of necessity and content itself with what
inevitably belongs to the understanding of \emph{given}
% --- p. 54 ---
\setcounter{footnote}{0}%
\opage{54}facts and scientific grounding of \emph{known}
laws. As in a holy twilight of half-transparent
necessity, the philosophy of nature can everywhere see the glow
of the divine freedom; but the will itself in
the inner infinity of the possibilities it cannot see: the
sun of freedom never rises in these regions of knowledge, it
keeps to the rim of the horizon.

The last proposition of the philosophy of nature is the first of religion:
nature is borne by the spirit, the laws are grounded in the will.
Now the sun of freedom rises, day dawns over the possibilities.
It is a sphere of cognition by itself, the eternal circle of light of faith and
of revelation. What do we see here in this new light?
The will in the omnipotent self, freedom in all its possibilities:
\emph{the miracle}! The miracle is that omnipotence is will, and
will omnipotence; the miracle is that the necessity of the whole world
has its original ground in divine freedom, the miracle
itself is the secret of wisdom.

The miracle of which we here speak is not phenomenon,
but being, it is a fundamental miracle. It has its unfailing
mark of certainty, not in outward signs, but in inward ones,
in our personal self, in the self-understanding of faith and will.
It is not first to be confirmed by lengthy hearings of witnesses
concerning outward miracles; it is, conversely, the outward miracles that,
when the critique of faith applies its measure, must all be traced
back to, and confirmed by, the ever-existing, eternally abiding,
never-changing fundamental miracle.

When we now, under these presuppositions, repeat that
the spirit is will, and its revelation % print: Aabenbarlse (misprint)
a miracle: then it is
surely evident that there lies hidden here not the least
afterthought or side-thought of breach of law, loose arbitrariness or
divine caprice. The spirit is being, and the spirit is
will. That the spirit reveals itself means, then: it shows
itself in the mirror of faith as it is according to the eternal
being of its will. What, then, is the spirit according to its revealed being?
Love, holy, unchangeable love. In this lies
that it is always ready --- for love is always
% --- p. 55 ---
\setcounter{footnote}{0}%
\opage{55}ready ---
that it will always come, when and where it, on
the terms of freedom, can come.

We know the spirit of truth by the Word of life. If we ask
why it then tarried so long before it came ---
with the Word of life to full revelation, then it is faith's
simple answer: it could not come sooner, it was hindered.
It was in the Word, and the Word was the light of the world, and
the light shone --- but in dark myths, and darkness and
myths comprehended it not. If we ask why it does not
come to us in tokens of power with the force of the miracle, as
it came to the fathers, then it is faith's simple answer:
it does not come with the force of the miracle, because it cannot,
it is hindered. It is the human will which,
on the way of freedom, has cast down the many stones and heaped up
the many obstacles.

``But'', one says, ``here, then, is a contradiction after all,
an intolerable contradiction, inseparable from miracle and will.
One of two things: either the spirit is omnipotent, or it is impotent.
If the spirit is omnipotent, what can stop its
coming? It comes and goes as it wills, when it
wills, and whither it wills; it comes with victory, shining,
warming: hearts melt at words of fire, like the ice
of the rocks at the glow of the sun. Thus an Omnipotent One comes. If
the spirit cannot come thus; if there is something that can
stand in the way, something real that can paralyze
its power, force it to stand still or at least to turn
aside and slink forward by detours: where then is the omnipotence of the will,
and what is the miracle?'' Undeniably! here is a contradiction.
God's will is omnipotent, man's will is impotent.
``And this impotent thing is then supposed to be able
to set itself up against the Omnipotent, this selfish element
in man's self-will is supposed really to have the power to
offer resistance to the spirit-will? It must be a delusion,
a deception.'' We answer: Yes! it is quite certainly a
delusion that wills of the day, even if they banded together
by the million into revolt, should be able to set limits to the eternal
will; it is a deception, a horrifying,
% --- p. 56 ---
\setcounter{footnote}{0}%
\opage{56}spirit-consuming self-deception. But this self-deception is a
real something, it is in man's self-will so real
that it can ``grieve the Spirit''.

If it depended on power alone, the spirit would be sure
of its victory: the slightest breath, and all the world's obstacles
would give way like fluff. But what would then become of man's
freedom? If the spirit were to use power and take
the refractory self-will by force, it would have to found
freedom by making freedom nothing: that is a
contradiction so clear that we see in it the impossible. If
the talk were of making the deaf hear, the blind see,
of opening the graves and raising the dead, then we could not
say: that is impossible for the will of omnipotence; for our knowledge does
not suffice. We can say that such deeds cannot be
explained according to the laws known to us; but we cannot
say that they stand in conflict with the laws, for the number of the as yet undiscovered
laws is infinitely greater than the number of the
discovered. All the contradictions that a frivolous knowledge in
its rashness has formulated in this respect are built
on sand and fall away of themselves when they are tested. But
the contradiction, that the omnipotent will should be able not merely
to compel the impotent --- there is no contradiction in that
--- but to compel it not to let itself be compelled, to compel it to
determine itself with freedom, that is a contradiction so clear
and evident that he who ascribes such a possibility to omnipotence
does not know what he himself is saying. The knot of freedom
will not be hewn through, it will be untied.

The untying of the knot is the dissolution of self-deception. Self-deception
is in man twofold, indeed threefold. It is in
thought and is called blindness; it is in the will and is called
defiance; it is in the feeling's unity of thought and will, and
is called hardening. It is difficult, if not impossible,
to get at the deception under any single one of these
forms by itself; the three fellow-interested parties, thought, will and
feeling, hold together in the common deception, they are like
conspirators: the one will not betray the other.
The understanding cannot be cleared, nor the blindness lifted,
% --- p. 57 ---
\setcounter{footnote}{0}%
\opage{57}so long as the will is perverse and the stream of feelings runs
against the truth. Feeling receives false impressions so long as
the understanding is clouded, and the will, infatuated by false conceptions,
keeps the entrance barred against every ray of light.
And would a will with false thoughts and
false conceptions, false images and selfish impulses
be able, on its own, to begin with the truth
out of itself and cut the deception through, as a
Chlotaldo hews asunder Morgana's pearl? The whole points
back to the unopened natural ground of the self, the unconscious unity
that bears thought, will and feeling in unexplained self-determining:
if the deception is to be able to be lifted, there must be in
this natural ground something that holds back and
keeps a possibility for the spirit.

Could we imagine self-deception continuing its course
through the three principal forms of thought, will and feeling
so long, so persistently and in such depth that
what is hidden in the ground of the self were brought up into the light of consciousness
and the last remnant of natural immediacy used up:
then self-deception would condense itself into a will of unconditional % print: fortætte (misprint)
denial of the light, a will darkened in spirit, a spirit-will
that was not human, but satanic. That under
this presupposition it would be neither thought nor feeling,
but the will itself, that in the very last would come to turn
the scale, is evident. Feeling is involuntary; the
feeling being does not have the impression in its power, and in this
there lies always an excuse for the deception of feeling. Also
the deception of thought has its excuse, so long as the unconscious
is still so predominant that a semblance can be kept up
that thought is continually about to seek
the true. Only when the denied truth stands
cleared before thought and the sting of pain goes through the heart,
does self-will have the conditions to transform self-deception
into absolute self-blinding and to drive hardening to the
utmost.

If such a self-contradiction were merely logical,
it would without further ado be referable to the impossible,
% --- p. 58 ---
\setcounter{footnote}{0}%
\opage{58}and the deception would glide past thought, like a shadow on the wall.
But a self-contradiction of this kind is more than logical,
it is personal. Already at the stage of the drive-will the
individual experiences that the true is by no means the
same as the agreeable, and the understanding on such
occasions often gets busy enough veiling the untrue
and adorning the half-true. At the stage of rational cognition
self-will can, in the name of fate, hold itself so
firm that it would rather be ``broken'' than ``bent''. If
now such a potentiated self-will comes into relation to the spirit
itself, the spirit of truth and of love, then it is by no means
given that it bends before the spirit, even if it
must acknowledge the spirit as a power. The opposite may happen.
The potentiated self-will feels that in all its impotence it
yet has a power, the power to defy the spirit. Seized by
the cramp of self-contradiction, it then seizes with the force of despair
upon the choice: rather to suffer all the torments of self-consumption
than to confess its guilt, bear the humiliation and bow
before the One mightier. When this choice has been made, then
the blindness is complete; then truth is lie, and lie is
truth, freedom bondage, and bondage freedom: the light of love
is extinguished in the darkness of hatred.

The same spirit of truth that awakens love
can also awaken hatred, for hatred is an inverted love, distorted in selfish
defiance. It is not the fault of love
that it awakens hatred; it cannot help it.
Love can bear everything for the sake of others, but self-contradiction
it cannot bear. It cannot deny
its being and say to itself: I must change myself, I
must see to becoming self-loving, that is to say: unloving, for
otherwise I awaken hatred. The being of love is: not to
conceal, but to reveal itself. When, then, the spirit of
the self-revealing love awakens hatred, the hatred
does not become slight because love is great; no, hatred feeds
on love, grows by love, strives to
measure itself with love and to become infinite, because
love is infinite. What is hidden in this, we shall in
% --- p. 59 ---
\setcounter{footnote}{0}%
\opage{59}this connection not try to bring forth, but only
hold fast that on the way of the will, in the name of freedom,
there really is a something that can grieve the Spirit.

We have considered self-deception from the strong side;
let us in conclusion reverse the standpoint and consider it
from the side of weakness.

There are indeed human beings, and their number is surely not small,
anxious souls, who sigh within themselves and say: ``We do not
resist the spirit; we feel that we need the spirit; we
pray for the spirit, but --- the spirit does not come.'' And there
are, by way of compensation, those who enjoy the spirit to excess,
the elect of faith, whose oily-browed thriving positively glistens
with the spirit: they chatter about the spirit and chatter with the spirit.
They possess the magic word in which the spirit dwells; so they know
the spirit well: they are, to put it plainly, comrades of
the spirit. There is in the relation between God and each
single human being a secret that only the All-knowing sees through; but if we look a little sharply at the human,
does it not show itself here too that self-deception is
twofold, indeed threefold.

There is a deception of conception which, forgetting
the omnipresence of the spirit, takes the word of revelation
in vain, fixes the spirit in an image and makes the image
outward. Under the image of an invisible, yet
finite person moving in space, ``the
Holy Spirit'' flies, up and down and down and up, a winged messenger
between earth and heaven. We take no offense at the figurative;
but when this figurative crystallizes in the fantasy
into a fixed reality, when the spirit at certain times
pays a visit to a human heart and goes away again with
a promise to return; when, we say, this image-being
gains such a power over consciousness that the one praying
sits and waits for the spirit as a guest to knock at
the door: dare we then wonder if such a waiting one can
come to wait in vain? To make the spirit of personality
into a mythical being, betrothed to Iris, who flies on
the rainbow, is to conceive it as an idol; but to conceive
the spirit as an idol is to grieve the Spirit.

% --- p. 60 ---
\setcounter{footnote}{0}%
\opage{60}The deception of conception just indicated joins itself with
a corresponding deception of feeling. % print: Følelsebedrag (misprint)
The delusion is that the
presence of the spirit is supposed to be the object of an immediate inward
sensation, that the soul whom the spirit visits
is supposed to be able to say to the hour when the visitation
took place, to be able to make a distinct separation
between two beings, its own being and the spirit,
to be able quite literally to converse with the spirit: ``I said
that to the spirit, and then the spirit said to me'', etc.
When this delusion gains the upper hand, and its hidden
influence is without doubt far greater than many
will rightly admit: dare we then find it strange that there are
sincere souls who admit to themselves that they have never
stood face to face with the spirit itself, never
at any time either seen or immediately perceived
God's own presence as Holy Spirit?

Then Holy Scripture must bear the blame. In opposition
to the image-believers, who, in order to take the word of Scripture quite
literally, make the supernatural unnatural, and set the
inward down upon an outward, there appear then ``the rational''
with an interpretation by knowledge, which disposes of the supernatural,
dilutes the word of revelation and has its distinctive
character effaced. Thus the book of the spirit becomes, by means of a
double self-deception, an equivocal book: the children in the nursery
keep to the pictures and play Christmas with the supernatural;
the grown-ups discover naturalness and fawn upon
mythology.

The deception of conception and the deception of feeling unite in a third,
the most dangerous form of them all, and that is the deception of the will.
The deluding thing, from which one proceeds, is the notion
that man's natural will can be left over as a
magnitude by itself beside the spirit-will. With some it is
conceived in this way, that the natural will, be it in
the form of drive-will or reason-will, is to be forced in
under the law of the spirit and whipped up to, if not precisely deserve
the spirit, then at least make itself worthy of the spirit,
work together with the spirit, and, if possible, by good
% --- p. 61 ---
\setcounter{footnote}{0}%
\opage{61}works acquire merits with the spirit. To this
belongs a peculiar form of the so-called strict, worried
Christianity, not the one truly, but the one sourly worried,
a Christianity whose love is a kind of
godly vexation, whose comfort is an assurance of the eternal
damnation of the glad, a Christianity that hates the world with
a yellow faith that has its seat in the gall. With others
man's natural will is quite rightly declared
unfit to do the works of the spirit. Everything rests without
further ado on the spirit and on faith, the new will comes of
itself. If the spirit wills, then there is will; if the spirit does not will,
well then, there is as yet no will --- the will grows
slowly. But that makes no difference: of faith there is plenty,
and spirit and faith are the whole. On these conditions
a fine reconciliation comes about between the man of spirit
and the natural man: they can now get on so well together,
for they both love freedom. Faith has its freedom
in the supernatural, and that cannot bother the
natural; the drive-will, which has its freedom in the natural,
can go about and graze wherever it likes, on the common pasture of
temporality, so long as it is otherwise well-behaved and
does not lay hands on faith. The reconciliation is so hearty that
it looks lifelike as though the two wills, that of faith and that of
the drive, had become one will. To this belongs a certain
formation of the so-called glad Christianity, not the one
truly glad --- for it is eternally certain to faith that the life
of the spirit is a life in gladness --- but the one gladdened
by naturalness, the coffin-glad, which finds it so lovely on earth,
because it is so sure of heaven.

When these self-deceptions are lifted, and they can be lifted
--- every human being has in the revealed word, the word whose
living, sacramental presence is to be found in the congregation,
the conditions for lifting them, when he earnestly wills --- then
the spirit comes.

In our present time various attempts are made
to reconcile the consciousness of the present with the existing religion.
One proceeds in this respect from the quite
% --- p. 62 ---
\setcounter{footnote}{0}%
\opage{62}certainly universally valid presupposition, that humanity has
a need of religion, so deep-reaching and essential a need
that states and peoples, civilization and morality, all civil
order and all the goods of a higher enlightenment would dissolve
into a chaos of raw passions and the features of humanity
be distorted into the ravenousness of the beast of prey and into the likeness of the ape,
if religion disappeared from human society.
One further holds fast to the in itself fully valid
truth that the religious cannot possibly be to us a living,
present thing, if it is not able in a spiritually powerful manner
to penetrate the consciousness of the present so completely
that we with full freedom live and act and move therein
as in our proper, spiritual element of life. It has quite
rightly been seen that it must be this spiritual element of life
that contains the first decisive conditions for the formation
of characters with a new outlook and a new will,
a will that with faith and courage dares to seize the tasks from
the ground up and boldly work at their solution. About all
this the more progressive in the age must surely all be
agreed; but, if we go closer into the matter, do we not
meet this strange thing, that the attempts begin by
overlooking the first necessary thing?

If the religious is to be put into force, then it is surely
altogether necessary that everyone begin by asking
himself wherein the religious essentially consists, what
religion originally is, and where the mark of truth is found,
the unfailing one, which can assure him that he does not
go astray and set out upon the uncertain. If the mark of truth
is reliable, then it must be found in his own
self, in the self-understanding of faith and spirit. This self-understanding
is an understanding of will, an understanding of the miracle;
for the miracle is religion itself. That a human will
stands in a relation of necessity to will-less things, in that
there is no miracle; that human wills relate to
one another as points of arbitrariness, all ruled by
overpowering laws, in that there is no miracle. But that God's
almighty will has entered into relation to man's
% --- p. 63 ---
\setcounter{footnote}{0}%
\opage{63}impotent will, has revealed itself as a human will,
and through this its self-revelation founded a relation of love,
a blessed reconciliation, that the two wills, God's and man's,
to all eternity can be as one will, that is in
spirit and truth a miracle. Those who think that they can best
put the religious into force by excluding the miracle
are like those who think that they can best draw the circle
by representing it as a square. To reject the miracle
is to reject religion; to accept religion as religion
is to acknowledge the miracle.

If the mark of truth stands firm in the inward, then there issues
from it the light that can illumine the corresponding outward.
The inward in faith is so far from being indifferent to
its outward that every believer must acknowledge that his self-understanding
in spirit is not owed to himself, but is awakened by
a word coming from without, and rests on the sacramental
miracle of the word, which with the omnipotence of the spirit-will communicates its
inward in, with and under a determinate outward. So here there is
thus by means of the miracle a continued, living reciprocal relation
between the inward and the outward. So long as all the outward
is conceived and judged as means and medium for
the inward, so long does self-understanding continue to enrich itself,
and the mark of truth holds good. Only when consciousness
makes an attempt to go around faith, behind faith,
in order to grasp the religious fact in pure, historical
externality, does the mark of truth fail, and groping and doubt
begin. It cannot be impressed enough that faith and
the content of faith are inseparable: every miracle is a caricature
when it is not seen with the eye of faith. Sacred
history unfolds for our contemplation a series of
pictures of life and leads us into a circle of facts
and phenomena in which the believers of the past, be they
eyewitnesses or not eyewitnesses, have seen and recognized the
tokens of omnipotence by which the saving, redeeming will of love
has revealed itself to them; the question is
whether in their faith we can recognize the essence of faith and share
faith with them. The understanding and appropriation of sacred
% --- p. 64 ---
\setcounter{footnote}{0}%
\opage{64}history is a continuing recognition of the life of faith,
its essence and content. Where this recognition
fails, there our mark of truth fails too, and the
handed-down does not touch us religiously.

It is a fully justified demand that the more recent age
has put forward, that research shall be free. On the conditions
here set up, the faith of revelation cannot have the
slightest thing to object to the unbounded
freedom of scientific research; for now it knows with itself that
no critical science, be its results never so negative,
is any longer in a position to undermine the religious foundation.
Should it even be proved scientifically that we had not
a single Gospel directly from any eyewitness, yes
even if it could be proved that our Gospel of Matthew
was from the second, and our Gospel of John from the
third century: all these proofs would still not
have the slightest influence on the determination of the content of
religion. Here it is not asked about the historical
intermediate links through which the communication of the truth has passed;
here it is asked about the truth, the revealed truth, and
the mark of truth, the unfailing one, is in our inward.

We live by the miracle and yet under the laws; we live
under the pressure of necessity and yet in freedom, in the midst of
the tangible facts of the material world and
yet in a present eternity, so truly have we the personal
certainty that will is spirit, and spirit will.

To measure one's faith by the warmth of feeling is uncertain,
for the measure of warmth is indeterminable, and the indeterminable
can deceive. To measure one's faith by the clarity of thought is
uncertain, for clarity is cold and the ideal in the mirror of thought
stands outside reality. But to measure one's faith by
spirit and will, by the spirit-will, which in the clarity of thought
has the warmth of inwardness and the reality of self-determination,
is certain: the measure of truth for spirit and faith is
a strong will.

% --- p. 65 ---
\setcounter{footnote}{0}%
\lecture{Sixth Lecture.}\opage{65}

A strong will wants to act: to choose, to decide, to do;
the drive-will wants this, the reason-will wants this, and in the highest sense
the spirit-will wants this. What, then, does the spirit-will choose? It
is divine in essence, and chooses the human; it
is human in active revelation, and chooses the divine:
in its choice it is divine-human. It is
in the Word as God that the Eternal himself elects every human being,
and it is again by means of the Word that man, by the power
of the Spirit, can make the election actual. Only when the
human will is at one and the same time chosen and chooses, is chosen
by God to choose God, to choose its own source: then
the self-deception of self-will bursts. The self chooses itself according to its
eternal destination, chooses freedom in unconditional dependence,
chooses to will One thing: to be the organ of the Spirit, chooses to will
God's will. The choice is not between This or That
among many goods; the choice lies within the will itself, between
a lower and a higher, a corruptible and an incorruptible will. It
is the deepest crisis of the life of will: the old will dies, in order to be
born anew; the drive-will is buried, in order to rise again transfigured
as spirit-will. Such a choice between life and death
is the choice between faith and unbelief.

Faith, it cannot be repeated enough, is more than
an emotion, more than a credulous acceptance of a prescribed
something that is to be believed; faith is absorption
% --- p. 66 ---
\setcounter{footnote}{0}%
\opage{66}in spirit, is the full surrender of the heart to the Word of life; this
surrender is more than an insight, more than a thinking,
it is an act of will, an inner deed, the first,
original action of the spirit-will. ``Give me thy heart!'' that is to say:
give me thy will.

The strong will always has a goal; the drive-will has one,
the reason-will has one, and in the very highest sense the
spirit-will has one. What goal, then, has the spirit-will set?
It wills the temporal, but for the sake of the eternal; the eternal
is the goal, the temporal is the means. The goal is: the personality's
temporally conditioned, yet strongly deep, many-sidedly actual
unfolding of the inexhaustible content of an eternal life.
Therefore this goal is both temporal and eternal.

The drive-will has its goal in the temporal; in all
naturalness it wills something present; the reason-will has its goal in
a present eternal, whose laws drain the lifeblood of the natural and
devour the temporal. The drive-will wills the human immediately,
but makes the personal temporal, and becomes
inhuman. The reason-will, in the name of universal necessity,
wills the human, but dissolves the self and thereby
becomes inhuman. The inhuman that lies hidden
at the bottom, under the reason-will as well as under the drive-will,
stems from the fact that both will so exclusively the
human, only the human. The spirit-will, on the other hand,
has in the unity of the goal a double goal; in the personal
sense it wills the human, because personally it wills the
divine. It does not will a divine for itself and
a human for itself; in the holy name of true personality
it wills the divine-human. Thus its
goal embraces both earth and heaven, both a present and a
coming, both a Here and a Hereafter: the divine-human unity
of the two-sided goal is life itself, the
eternal life.

That eternal life is life itself can be seen from the laws.
For it is not the case, as many seem
to assume, that there should be, by origin, two kinds of
life, one that springs from temporality, alongside an
% --- p. 67 ---
\setcounter{footnote}{0}%
\opage{67}other, which is from the Eternal. Those who cling to
the notion of a life whose origin, development, and
conclusion should belong exclusively to temporality have
not duly accounted to themselves for what temporality is, and
what is understood by that which we call life.

By no means is everything temporal that takes place in time;
for every development, including that of an eternal content, must take place
in time: time belongs to eternity. Temporal, on the other hand,
is that which takes place in time in such a way that its
activity and being presuppose parcelling-out and are conditioned by
composition. Changeability, mutability, corruptibility,
the progress of the one out of the other and its dissolution
into the other, thus parcelled-out otherness, is the mark
of temporality. In this sense all material
formations are temporal.

It is otherwise with the Eternal, whose essence is in all things like
itself. The all-life is eternal, a self-being, which through
matter, the temporal other, develops itself, only itself
and its content. The life in everything living begins from within and
maintains itself, in the working-over, reshaping, and
appropriation of matter, in its inner circuit; in this circuit
the Eternal shows itself. There is declamation about the eternity of matter, about the
eternal stars, the eternal mountains, and so on; it is mere declamation.
The existence of matter, of the stars, of the mountains is
a selfless existence, an existence in otherness, thus in
temporality. The billions times trillions of years with which
people so readily toss about are not able to make all this
temporality into a real eternity. But the all-life, which
in the midst of otherness defeats parcelling-out, continues and
preserves itself, is eternal. There is more eternity in the cells of the plant and in
the flowers of the field than in all ``the eternal
stars'' together. The phenomena of life are corruptible, but
the essence of life in its quiet circuit is incorruptible; it is not
composed of an other and is not dissolved into an
other: it remains in itself, unchangeable in its changes,
imperishable in its interior, although the exterior
perishes.

% --- p. 68 ---
\setcounter{footnote}{0}%
\opage{68}Let a man reflect properly on the difference between the living
and the lifeless, and he will convince himself that
every being we regard as living, even if it lives
only in a single cell, must be regarded as a self-being,
whereas matter, the material, with its forces and laws,
is only a being of otherness, is lifeless, because it is
selfless. Herein is seen a thoroughgoing difference between physical
laws and laws of life. The physical laws are laws of temporality;
only the laws of self-development are true laws of eternity. From mere
parts of matter and physical laws the origin of life cannot be derived:
to derive life from the lifeless, the self-being from the
selfless is just as impossible as to derive something from nothing.
Only when the temporal comes under the eternal, and the
physical laws are taken up into laws of self-development, do
we recognize the laws of life.

On what, then, does the separation between temporal and
eternal life rest? Not on the essence of the all-life, for that is from eternity
a self-being, but on the particular appearing of life in the
living individuals. In this respect the difference between
plant life, animal life, and human life is of thoroughgoing significance.
In plant life the self-being is so fluid that an
enduring self does not form; in animal life the self-being receives
a natural self, for the animal has a soul; but
only at the stage of human life is the germ laid for the development
of a self that is conscious of itself, a spiritual, personal self.

The question how the living individuals can
continue their self-development beyond the bounds of temporality
no science can answer. Such an infinitely
continued development of the individual life is, from the standpoint of
natural experience, improbable, but, according to the
laws of life, nevertheless a possibility. It is this possibility
that faith seizes, in that it discovers the eternity of the self, dissolves
the self-will ensnared in temporality, and transforms it into
spirit-will. The full assurance of the eternal life
of individuals is conditioned by spirit-will.

It is a frivolous notion that in this present life there
should be no personal certainty of an
% --- p. 69 ---
\setcounter{footnote}{0}%
\opage{69}eternal life, but only after death. Should we discover only after
death that there is an eternal life, the discovery might
become terrible.

Self-will, whether it otherwise appears in
the form of drive-will or as reason-will, denies the
eternal life, and breaks its laws. The drive-will asserts
the temporal claims of individuality, but excludes the eternal
ones; that is a breach of the law, a breach of the law of life. The reason-will
asserts claims of eternity, but in a selfishly selfless
manner, to the destruction of individuality; that is a
breach of the law, a breach of the law of life. Such a self-will
can make a fuss about necessity, bow and scrape before the
immutability of the laws of nature as much as it likes, it nevertheless
breaks an immutable law, the holy law of freedom, of personality,
the fundamental law of the eternal life. But to
go through death into eternity as a known
criminal against the fundamental law of the eternal life is, after all, a
precarious matter.

At the stage of the spirit-will the relation between the
temporal and the eternal presents itself as a relation of discord and reconciliation.
Only in so far as the temporal is in every respect
made a means for the eternal can the spirit-will use the
temporal and, despite resistance, hindrances, and adversities,
be reconciled with it. Thus there may be striving and suffering in the
temporal, but when there is striving and suffering for the sake of the eternal,
it is recognized that strife and suffering furnish conditions
for developing the strength of the will and heightening its force; with
this recognition the assurance of eternal life grows: the
Eternal comes constantly nearer, so there is a reconciliation.
But although the life in this reconciliation is \emph{eternal},
yet it is still not \emph{blessed}. Blessedness is only in
bright moments; the discord is not removed: the opposition
remains between the blessed and the not blessed. That is the law of the
life of will within the bounds of temporality.

It is otherwise when temporality is exchanged for eternity,
when \emph{Here} is swapped for \emph{Hereafter}. In eternity
itself the two beings, the temporal and the eternal, are enemies
% --- p. 70 ---
\setcounter{footnote}{0}%
\opage{70}to the death, irreconcilable enemies. In the kingdom of eternity
no temporality is tolerated, for here no corruptibility is tolerated;
it follows from the nature of the case, it lies in the law of eternity.
All outward temporality, all marks of natural
corruptibility death can take away; but the corruptible
that is in the inward, when the self with its self-will
is in corruptibility, in spite of all its bustling activity yet
in its self-deception sheer finitude, sheer temporality, sheer
corruptibility, this mark death cannot take away;
for death has no power to blot out the personal in
the being of man.

So this is the unbreakable law of the eternal life:
``\emph{Either --- Or}''. Either man has his innermost
personal life in the spirit-will, and then the spring of eternity
comes and unfolds life into blossom; that is the blessed. Or
man has his life in his self-will, has in the universal of
self-will willed himself and his dissolution. Then
this self comes, with its self-deception, as the enemy of eternity, into
eternity. It wants to be dissolved, and yet cannot be dissolved,
for the eternal in the self asserts its right. It has
not willed an eternal life, so it must suffer all the torments of
self-contradiction; it has not willed the will of all-love,
so it meets the Power, the omnipotence of the judging will; that
is the unblessed. The opposition of blessed and
not-blessed in temporality can thus in a Hereafter transform itself into
an opposition of eternity between the blessed and the unblessed; that is a
judgment according to the law.

``But,'' say the spokesmen of enlightened reason, ``this
unblessedness is precisely the inhumane. If the spirit-will believes it
can cow us with unblessedness, to put it plainly, with the possibility
of an eternal damnation, then we have here a
palpable proof of what it is to which the whole
points in the end. An eternal damnation! It is a new
rehash of the old contradiction, a contradiction in concepts,
which science has long since resolved, --- the unblessed, the
barbaric, the inhuman!'' Our answer is brief. As
life is, so are its laws, and the laws are indeed ---
% --- p. 71 ---
\setcounter{footnote}{0}%
\opage{71}``unchangeable.'' If the fundamental law of freedom, of personality,
is barbaric, then the eternal life is certainly barbaric;
but then all talk of inner independence, personal
goal, struggle for freedom, and whatever else one likes,
when one speaks in earnest, is meaningless, barbaric talk as well.

Let us keep to what we know, the laws of the eternal life
within temporality, and by way of conclusion cast a glance
at the relation between \emph{the religious} and \emph{the humane}.

As regards this relation, we dare not
stop at the external and call history to witness
how often religion has wanted one way, and humanity
another. The question is: what does religion demand,
and what does humanity demand? Can the religious and humane
demands be united or not? Nothing can be more
enfeebling to the spirit than the notion, conceived in a foggy piety,
that life in God and life in the world should be
incompatible things, that a vigorous, capable, active life in the
real world should be a sin. When those who
call themselves ``God's children'' are nevertheless compelled to
``work and look to the profit'' just as much as the ``children
of the world,'' they console themselves with ``hard necessity''
and ``God's grace,'' and drift along in a pious bit of petty twaddle about
``the forgiveness of sins,'' which does not further the humane, and
has no warrant in the religious. For what is it
that religion originally calls sin? It is the state
in which man finds himself outside of spirit, the
state in which the individual has fallen prey to the drive-will and
the self-will forsaken by God. Man, says religion,
is born in sin; for he is born under the drive-will. From
the state of sin spring the so-called actual sins;
they pass through feeling, imagination, and representation and receive
the stamp of individual imputability, in that they are taken up
and carried out by self-will.

If we consider these actual sins more closely, it
appears that they disturb life in the world just as much as life
in God, that they are equally destructive of the humane and
of the religious. It is a peculiar kind of pious naïveté,
% --- p. 72 ---
\setcounter{footnote}{0}%
\opage{72}when a soul-shaker, who really means to go for the world's
ungodliness, sets about thundering against murder, theft,
excess, avarice, malice, deceit, licentiousness,
and so on, as clear proofs of rebellion against God, of transgressions
of his strict commandments. Anxious consciences get
an impression that, if the Almighty had not so literally
forbidden all such things in his holy law and threatened
the transgressors with hell, then such gross outbreaks
of human self-will would not properly speaking have been
sins. It is forgotten that the commandments of omnipotence are commandments of love,
that evil is forbidden because it is in itself evil,
and is evil because it undermines the human. A
God whose will demands something other of man than
what every human being who in faith recognizes his eternal
destination must demand of himself is an idol; in
other words: a religion that stands in conflict with true
humanity is just as false as a humanity that stands
in conflict with true religion.

This insight is most closely connected
with the right understanding of what forgiveness of sins is.
It is a childish notion that the Father in heaven, who
has become angry because his strict commandment has been transgressed, can,
when the Intercessor steps between, be mollified and become good
again, when the sinner creeps to the Cross and humbly begs
for forgiveness. The childlike in such a notion we ought
to honor and respect, for it is eternally certain that the relation to God
must on man's side remain a childlike relation.
But the childlike is essentially different from the childish.
And the relation becomes childish among the
``sweet'' Christians as well as among the ``sour,'' when the eternal
fundamental law of forgiveness of sins is misunderstood.

The Cheerful, who are so rosy-cheeked because in the midst of
``worldly twaddle'' they can do more or less what they will,
and yet always be assured of ``the forgiveness of sins,''
because it stands in the third article of the faith, are not
childlike, but childish. Instead of measuring sin
by the standard of self-contradiction, as humane as it is religious,
% --- p. 73 ---
\setcounter{footnote}{0}%
\opage{73}and of knowing the seriousness of repentance and the reality of forgiveness of sins
by the renewed tensile force of the reawakened spirit-will,
they set the Law aside, and wax and grow in the
conceit that an old psalm is quite as good as
a new will.

The sour Christians, who hold to the laws of compulsion, have
trouble enough, not only with their own, but also,
and chiefly, with other people's sins: no wonder
that they are so pale. When the great register of sins with
adultery and murder and luxury and excess and deceit and
theft and so on is unrolled, it escapes their notice
that the soul-judge stands with one foot in heaven,
but with the other in the town hall; it escapes them
that the great sins on the reckoning-board are at bottom just as
inhumane as irreligious, that the police are constantly on the
trail of the grossest, although the police are after all not a
religious, but a plain and simple human institution: this
grand zeal against ``the sinners'' and ``the sinners'' and ``the sinners''
is, when seen in its relation to spirit, less
serious than childish.

The more religious knowledge clears itself, the more
distinctly it will be seen that life in the world must be an
essential condition for life in God, as surely as the religious
life in the highest, fullest measure is a life of will. One thing is \emph{needful};
but for the sake of this one thing a worldly content
of feelings, representations, interests, and labors is needful.
Sin can now take two forms. Either the self lacks
the one thing needful, love of spirit with faith and hope
is left out, and self-will greedily throws itself upon the
manifold: that is sin in the form of worldliness;
or the self expends all its force of will on separating itself
from the real worldly content and burning up in God,
whether by the way of finitude through the good works of
arbitrary renunciations, or by the way of infinity by
mystically plunging itself into the sea of light of the unnameable all-love:
that is sin in the form of misunderstood godliness.
Life in God means that man unconditionally
% --- p. 74 ---
\setcounter{footnote}{0}%
\opage{74}has given up his self-will and recognized his absolute dependence
on God's will. Life in the world means that
absolute dependence grounds a real freedom,
that the spirit-will makes man independent. The Spirit
proves precisely the force of its personality in this, that man
with strong will and a sure hand takes hold of the labors
of real life.

The naive separation of the Middle Ages between \emph{godly}
and \emph{worldly} works can stand neither before the judgment of religion
nor before that of humanity. All moral actions
are equally religious and humane, as surely as they
proceed from the spirit-will.

No work in the world that stands in any relation
to the personal life is too small for the Spirit; and no
external finite work, even if it outdid the greatest
in a Loyola's heaven-storming asceticism, is great enough
to express the Spirit as Spirit.

There is in all works, be they ``godly'' or
``worldly,'' sin, when the spirit-will is lacking; but in all
works, be they called worldly or godly, the sinful
is sublated, in so far as the spirit-will takes effect.

There are not three kinds of works: some that we are to
do for God's sake, others that we are to do for our neighbor's
sake, besides those that we are to do for our own sake.
Love God above all things and thy neighbor as thyself!
These great commandments of the Law are one commandment. It is
impossible to love oneself in truth without loving God,
and it is impossible to love God without loving one's neighbor.
The religious is that love to God stands highest; the
humane is that love to ourselves stands inseparable
from love to the neighbor.

The old unfruitful school question, how we are
to find out with certainty what God's will is with
regard to the single works of life in the various
cases that meet us, falls away of itself from this standpoint.
Existence sets man an unsurveyable multitude
of highly different tasks. The solution of none of these tasks
% --- p. 75 ---
\setcounter{footnote}{0}%
\opage{75}is left to man's selfish discretion.
The religious demands that they all be referred, immediately or
mediately, --- in faith --- to God's will. But
none of them all is, in its immediate singleness, literally
determined by a command dictated by God; they are to
be grasped all with the freedom of the Spirit. The humane comes into
its right in that the human will as spirit-will is called upon, in
each case that meets it, to act out of the
interior of personality.

The fundamental relation between the religious and the humane
can, according to what has gone before, be determined exactly by
the relation between drive-will, reason-will, and spirit-will.
The reason-will, as we have seen, signifies a stage of transition,
and does not let itself be carried through; it hovers constantly
between the religious and the irreligious, between the humane and the inhumane.
The drive-will, as it begins to become rational, goes without
further ado for the humane. Even if it is influenced
by spirit, it does not therefore give itself up, but makes
a deal with the religious, and seeks in the best way to
draw advantage from religion in order to preserve and further humanity.
To this corresponds the often somewhat devious
use that politicizing societies, from considerations of prudence, have
become accustomed to make of the religious institutions.
The spirit-will, on the other hand, stands in a double relation to the
humane, a hostile and a conciliatory one. It stands in conflict
with drive-humanity, for it stands in conflict with the
natural drive-will. It places itself decisively on the side of
religion, directs its attacks against the humanity that denies spirit,
tears down its altars and burns its idols.
Thus it looks as if it were barbaric. But these
attacks on false humanity are in principle for the benefit of
the truly humane; precisely in that the humane
unconditionally comes under the religious, it is born again
to a higher life: the spirit-will intervenes in the worldly,
plants its banner of freedom and goes into battle for the right of
humanity.

% --- p. 76 ---
\setcounter{footnote}{0}%
\opage{76}The two-sided relation between the spirit-will and humanity
can be demonstrated in all the circles of society, in
the family, in the civil sphere, in the common
purposes for which people work with ideas.

Marriage is the moral foundation of family life.
Whence does it spring? what validity has it? is it
dissoluble or indissoluble? Those who live with the unpainted
drive-humanity must naturally declare marriage
a natural product of the drive-will. The humane is
then that marriage is conditioned by mutual love;
the immediate surrender of immediate hearts, with sacrifice
and accidents, is the happiness of marriage. If it now happens
that the fire of love goes out, and it can so easily happen,
since it is after all drive-love, then marriage must be dissolved,
in order not to become unhappy. If individuals do not
simply have the right to dissolve it themselves, then it is naturally a brutal
inhumanity of a slavish establishment, which dares to
set barriers to free love in ``the free
flesh.'' That the spirit-will cannot reconcile itself with such a
humanity is comprehensible. But the society that has
stopped at getting drive-humanity rouged up with
a few moral beauty-spots, without wanting to enter
into the unconditional demand of religion, does not satisfy
the spirit-will either. From this standpoint one arrives at the result
that marriage is certainly, according to the divine law,
indissoluble, but yet, when special dispensation is obtained,
be it from the Pope or the Chancery, nevertheless dissoluble
out of regard for the humane. Against such a half-religious,
shabby humanity earnest preachers of penance
have at various times raised their commanding voice: ``the
divine commandment must be heard, recognized and obeyed, even if
the humanity of the whole world is to perish by it!'' Thus
it looks as if religion were cruel and despised
the humane. And yet this is a monstrous misunderstanding.

The religious wants to rule absolutely, unconditionally; but not
by external commandments of law. The external commandment of law, whether
it is built on passages of Scripture or on chancery rescripts or
% --- p. 77 ---
\setcounter{footnote}{0}%
\opage{77}on both, must by its nature always keep to
finite conditions, and in its externality can
never make claim to hold as absolute. The
Absolute lies in this, that the divine commandment in the name of the Spirit
is a commandment of love, a commandment of personality, of independence,
and of freedom. The religious understanding
is that the commandment, in all its severity, does not wish to impose on individuals
external compulsion, but to convince them of the
claim of the spirit-will and to catch them in their own conscience. If
a religious individuality, after earnest self-examination before God,
comes to the humbling certainty that the marriage entered into has,
for whatever reasons, become an empty semblance, a yoke of bondage,
a snare, then such a personality, be
it woman or man, dares to step forth independently, to break its
yoke and go its own way in spite of the busyness of tongues and
the well-known Pharisaism of a righteous neighborhood: that courage
faith shall give. That is the humanity of religion, and
truer humanity cannot be given.

The moral foundation of civil society is the right
to property; where the right of property is not respected, no
social order can subsist. The enterprising citizen gains influence
in proportion as he extends his possessions.
The civil juridical personality appears in the form
of property: the greater the property, the more personality.
That is natural humanity. The drive-will has its most important
account to settle in measures of Mine and Thine. Thus
a human society forms with human purposes. If
the religious is now called to help to consecrate and bless all
this humanity, then personal property of course gets
high thoughts of itself. Property becomes capital,
and the power of capital is irresistible. Capital teems
under the blessing; it grows and multiplies, makes itself
the earth's master and rules over the fish in the sea and
over Majesty on the throne, over the laborer in the mine
and over every unpropertied beast that creeps upon the earth.
In order to secure all the fruits of the blessing, drive-humanity
becomes godly: the rich man goes to church on
% --- p. 78 ---
\setcounter{footnote}{0}%
\opage{78}Sunday, gives sums to the public cause, and so on; at the
standpoint of the drive-will all of it honorable, handsome, and
humane. But what does religion say? Away with all of it!
``Sell all thy goods and give them to the poor;''
``it is easier for a camel to go through the eye of a needle
than for a rich man to enter into the kingdom of God!'' --- How offensive.
Do you hear, say the rational, how fanatically
one-sided religion is, how little it understands
what is in reality humane!

To take away the offense and calm minds
there comes a softly sweet soul-comforter with the interpretation, as scientific
as it is humane, that it is by no means
religion's meaning that we should literally give away all
our goods: it only wants to awaken our conscience and
put a fright into the avaricious. The right meaning is
that we should think of the poor, give something away of our goods,
a good part, especially for godly use, but not all; the
most we should keep ourselves, so that we may truly learn to manage
it for the benefit and joy of ourselves and our heirs.
But what does religion say? Away with the whole! ``Sell
what thou hast, and give it to the poor.'' That is the way
to perfection, the way of renunciation to ``the kingdom
of heaven.'' Thus the commandment was understood in those days, when
the hermits of the desert and the strict monks of the Middle Ages had not yet
learned to tamper with the religious-humane.

Yet, one objects, ``the great progress that the present
has made consists precisely in its wanting to reconcile religion
with humanity.'' Indeed! But the reconciliation surely must
not be grounded on a sly deceit. The misunderstanding
of those ancients was by no means that they took the religious
demand to be unconditional. In the literal understanding
a holy Anthony, a Benedict of Nursia,
a Francis of Assisi stand infinitely higher than all the
mediation-ethicists of modernized Christianity put together.

The misunderstanding of the ascetics was by no means that
they took the commandment of renunciation too strictly, but on the contrary that
they did not take it strictly enough. They conceived the command of
% --- p. 79 ---
\setcounter{footnote}{0}%
\opage{79}renunciation of the world as an external command, a compulsory command
arbitrarily prescribed for the attainment of perfection. Granted,
this commandment was to them also a commandment of the heart, a commandment of love
toward deliverance and blessedness; but the commandment of
love was not conceived in the spirit of love, for it
was not conceived in the spirit of personal freedom. Thus
unfreedom triumphed, and within religion there arose two
kinds of morality, a higher for the perfect and a lower
for the imperfect. It was this blunder that made
the medieval religiosity inhuman.

Only when the commandment of renunciation is appropriated in faith
as an inward commandment, a commandment of freedom, of love, and of
personality in all its severity, does severity
again become mildness, and the humane is born again in the religious.
It looks so severe that man is to part in the external
sense with all his goods; the drive-will sees no
back door through which it can slip away from so pressing
a demand. But when the external is converted
into the internal; when the commandment turns and says: keep
what thou hast in the external; but let it go in the internal,
be with all thy property as one who owns not a farthing,
then, one thinks, religion is very mild: the drive-will spies
the back door, smiles in its beard, and has its own thoughts. And yet
it is so that the holy commandment of love, understood in spirit
and in truth as a commandment of freedom to personal self-development,
is in all its mildness far severer than any
ascetic commandment of compulsion. It is the force of the spirit-will
that is to be tested in humanity.

If the religious principle for true humanity
is not acknowledged, the strife over Mine and Thine will develop
into a barbaric strife, and in a barbaric strife victory
will in the end be on the side of the elemental forces. But if
the principle of humanity is fully acknowledged, no less
on the side of the Possessing than on that of the Non-Possessing,
the world-struggle will go its course, but humanity,
supported by the force of the spirit-will, will conquer.

% --- p. 80 ---
\setcounter{footnote}{0}%
\opage{80}If we cast a glance at the different labors that
are performed in the service of ideas and are conditioned by special
gifts, there shows itself an opposition between idea and
spirit. How great this opposition is can be seen in
Greeks and Jews. The Greeks have an overflowing
wealth of fruitful ideas, but what we in the proper sense
must call spirit they do not have. The Jews stand
at the stage of the Spirit under the will of the law-giving Jehovah,
but have a striking lack of productive ideas.
The Jews represent the religious, the Greeks the humane.
One will on that account find it natural that Christianity
in its great struggle of entry into the world began by
breaking the staff over all the works of genius of paganism; it
was in the eyes of the educated an eloquent proof of fanatical
inhumanity. One will further find it natural that
the humanity which an age that denies religion wants to
found on the modern dogma of the absolute unity of drive-life and
reason-life turns away with disdain from
``the Semitic race's'' stiff-minded narrowness
and seeks its models under the Greek sky. One will
finally find it in order that those who here dare not
see the hostile in the relation between a revealed religion
and a pagan humanity, because they dare not admit
to themselves that two absolutely heterogeneous principles must fight
to the death until either the one or the other gains
the power, scrape through by a mediation after their
own heads.

On the stage of half-religiosity, as half-human as it is unprincipled,
the stage of mediation, we then meet again the medieval
separation between the godly and the worldly, but with a
modernized cut. ``All the ideas of humanity are justified,
of course; talent and genius are good gifts,
and all good gifts are to be used.'' But now there arises the
difficulty that by far the most human ideas are by nature
pure heathens, and it will hardly do in our time
to baptize ideas, as in the Middle Ages they baptized bells.
One must then make a due distinction between those on the
% --- p. 81 ---
\setcounter{footnote}{0}%
\opage{81}one side that cannot be brought under the religious ---
physics, for example, and comedy --- and those on the other
side that can be used partly immediately, partly mediately
in the service of religion. The first kind one must let
go their own way in all worldliness and only take care
that they keep within certain bounds. The latter
kind one must prize, for they are mediators between
the religious and the humane, according as they are used
more or less godly. It should then, as regards the application
of talent and genius, according to this view,
be more godly to build churches than to build town halls,
more godly to play the organ than to play the violin,
more godly to paint saints than to paint cows. If
one nevertheless does not openly want to acknowledge such
a view, it is presumably because one is loath to acknowledge
the confusion that lies at the bottom of
the view; for it is precisely the confusion of what
concerns the idea and what rests on the will
that muddles the conception of the religious as well as of
the humane.

The content of the form of the idea is always from one side or another
a worldly content; the idea of God itself makes no
exception to this. --- No idea is in and of itself either godly
or ungodly; in this lies that all ideas without exception
can be brought under the spirit-will, in so far as the spirit-will
by no means wants to cow, but on the contrary to nourish, purify, and inspirit
the drive-will that gives birth to them into the world. Every idea
has a goal in itself, and every idea is a means for
the Spirit. To deny the development of ideas its natural right
is inhuman; but to set the absolute goal in any
idea whatever and deify the idea, whether it be
an idea of science or of art, the idea of a system of state
or the idea of a system of church, is irreligious.

Where the spirit-will is lacking, the ideas become selfish;
the tension and discord that arises among them sets
the passions in motion and gives them the rein. A billowing
sea of selfishly ideal endeavors can be set in uproar;
% --- p. 82 ---
\setcounter{footnote}{0}%
\opage{82}the storm can rise to such a degree that a richly gifted
people goes under in the strife. We have in the Greeks a
tragic example.

Natural humanity cannot bear itself,
it must be borne by the Spirit. The many conflicts of ideas are
tasks for the will; they form for the will's sake, they
call it to exertion, they demand its decision.
In a constant struggle against selfish hindrances we have
the conditions for the personal development of the force of the
spirit-will.

In these fundamental laws of existence it can be seen in faith
that there must be an inward connection
between the present and the coming, between
Here and Hereafter. This connection cannot be grounded
by an external consideration of ``the resurrection of the dead.''
The great pains which Christian dogmaticians, especially the
older ones, have taken with a --- of course ``scientific'' ---
depiction of the coming glory, of how the
dead are now actually to rise, at what age, in what
condition, one-legged or two-legged, with wig or without
wig, is wasted trouble % print: Uleilighed (misprint)
The whole question points
back to inwardness, personality, the spirit-will,
freedom.

The manifold exists for the sake of the one, the
temporal exists for the sake of the eternal: the eternal life is an
inner life, which nourishes itself on recollection. Everything that
is experienced in the external leaves a trace in the internal,
that is to say: it is \emph{inwardized}, it is recollected. The
more strongly the will is engaged in the experience, the deeper is the
trace of recollection graven. That these traces could ever be effaced
is an imagination, a sensuous self-deception, which
stems from the organs becoming dulled. In the life of the soul, in the ground of personality
they are never effaced: the laws of life all point
to one and the same fundamental law, the law of recollection.
Much is said and much is heard, much is remembered and much
forgotten in this temporal, but the Spirit listens and keeps;
it forgets nothing: in eternity nothing is forgotten.
% --- p. 83 ---
\setcounter{footnote}{0}%
\opage{83}Memory can be weakened, the understanding clouded, and the
light of consciousness extinguished; but deeper than all this and
as the ground of all this the nature of the self lies preserved and
hidden in the Spirit's recollection. Out of the night of recollection
everything comes forth again, and day dawns anew: self-consciousness,
understanding, and memory awake with all that was experienced
and first and last with all that was earnestly willed.
What in this temporal has been lived and enjoyed, for which labor,
suffering, and strife have been spent here, what the will has attained, and
what it has inwardly striven for but not achieved, of all
this nothing is lost; in spite of sensuous forgetfulness nothing
is forgotten, it is all kept, buried and kept
in the lap of recollection for renewed review.

Our subject is not exhausted; it is in itself
inexhaustible. Far from the conceit of being able to exhaust
the inexhaustible, I have in these lectures confined
myself to illuminating some features of the great picture of the life of will,
in order thereby to make a contribution, however imperfect,
with the wish that everyone who will take his experiences of life
to help may earnestly become aware for himself of what
spirit-will is, and may clarify to himself the conditions for a strong
will.

\end{document}
